\section{Equality as a constrained minimum of the original data}
\label{sec:variation}

We now assume the hypotheses of the rigidity assertion, including equality,
and return to the original pair $(g,K)$.  The only conclusion needed from
the numerical argument, Theorem~\ref{thm:numerical-input}, is its inequality for all sufficiently small
admissible variations of these original data.  Those variations need not
satisfy the additional outermostness or area-minimization hypotheses of
the rigidity assertion.  No equality statement for a Riemannian
deformation, and no rigidity theorem in another dimension, is used here.
In this section $u,X$ denote the lapse and shift obtained as constraint multipliers.

Our goal is to show that the area multiplier is supported entirely on
the original boundary. We first retain all minimizing frontiers in a
relaxed variation formula, then remove their interior support before
invoking smooth outermostness.

Put
\begin{equation}\label{eq:variation:constants}
 r_h=(A/\omega)^{1/k},\qquad m=\tfrac12r_h^{k-1},\qquad
 b^0=E/m,\qquad b^i=-P_i/m,\qquad c=(k-1)/r_h.
\end{equation}
For nearby charges define $\mathcal E=b^0E_{\rm new}+b^iP_{{\rm new},i}$.
The reverse Cauchy--Schwarz inequality for future timelike vectors gives
$\mathcal E\ge (E_{\rm new}^2-|P_{\rm new}|^2)^{1/2}$.  At the original
data, $\mathcal E=m$, and, for $F(a)=\tfrac12(a/\omega)^{(k-1)/k}$,
\begin{equation}\label{eq:variation:area-derivative-constant}
 2k\omega F'(A)=c.
\end{equation}

\subsection{The relaxed area envelope}

Apply Section~\ref{q:geometry:section} to the original exterior
$(\Omega,g)$ and its boundary $S$. Its hypotheses follow from the
smoothness, completeness, compact-complement one-end condition and
asymptotic decay in Definition~\ref{def:data}. Use the compact
filling $O$ constructed there and write
$\mathfrak T=\mathcal T(g)$ and $\Gamma_C=T_C$.
The boundary is outer area-minimizing, so its area $A$ equals the
full-cut infimum. In particular, the original filled obstacle is
itself a minimizing competitor. The normalizations of the constraint
densities play no role in this geometric step.

Several enclosures may attain $A$. We retain all of them: under a
metric variation the right derivative of their least area is the
least of their first variations. This compact family will carry the
probability measure in the multiplier identity below.

\begin{lemma}[Envelope and local geometry]\label{lem:variation:envelope}
The minimum perimeter is $A$. The family $\mathfrak T$ is compact
in $L^1$, its members lie in one compact set, and convergence within
it is strict convergence of vector perimeter measures. The same
uniform confinement and subsequential convergence to $\mathfrak T$
hold for minimizers of the smooth metric paths used below.

Each $\Gamma_C$ is a multiplicity-one almost-minimizing boundary,
with singular set of Hausdorff dimension at most $n-8$. Its free
regular part is smooth and minimal. Every contact point with $S$ is
regular, with a $C^{1,\alpha}\cap W^{2,p}$ graph for every finite $p$.
Whenever the frontier is smooth, its closed exterior is connected,
complete, one-ended, and outer area-minimizing. Near a free regular
point there are uniform smaller cylinders in which all minimizing
frontiers meeting the cylinder are single graphs. Distinct graphs
are disjoint, and their ordered heights satisfy
\begin{equation}\label{eq:variation:graph-separation}
 \|h_1-h_2\|_{C^j(V')}
 \le C_j|h_1(y_0)-h_2(y_0)|,\qquad V'\Subset V,
\end{equation}
with constants uniform in the two members. For $g_t=g+th+o(t)$
among these paths,
\begin{equation}\label{eq:variation:envelope-derivative}
 \left.\frac{d}{dt}\right|_{0+}a(g_t)
 =\min_{C\in\mathfrak T}a'_C(h),\qquad
 a'_C(h)=\frac12\int_{\partial^*C}\tr_{T\partial^*C}h\,dA_g.
\end{equation}
\end{lemma}

\begin{proof}
Proposition~\ref{q:geometry:enclosure} gives the full-perimeter
realization and regularity with obstacle $S$, without invoking any
rigidity assertion. For the paths below, the compact variations vanish
in the tail, while the charge prototypes and strict conformal
direction have metric perturbation $O_2(r^{2-n})$. After decreasing
the parameter interval, these are uniformly asymptotically flat,
uniformly positive, and smoothly controlled on the compact region.
Thus Proposition~\ref{q:geometry:envelope} applies, including its
localized plane-measure continuity. Proposition~\ref{q:geometry:sheets}
gives the graph cylinders and Equation~\eqref{eq:variation:graph-separation}.
\end{proof}
\subsection{Separation with active constraints}

Let
\[
 \mathfrak A=\{(x,v):x\in\Omega,\ |v|_g\le1,\ \mu(x)+J_x(v)=0\}
\]
be the closed set of active rays.  For a variation $(h,p)=(g',K')$
transport the original ball isometrically by the metric square root;
thus $v'=-\tfrac12h^\sharp v$.  Denote the resulting derivative of
$2(\mu+J(v))$ by $\mathfrak C'(h,p;v)$.

\begin{proposition}[Constraint and area multipliers]
\label{prop:variation:multiplier}
There are a probability measure $\pi$ on $\mathfrak T$, a locally finite
nonnegative measure $\Lambda$ on $\mathfrak A$, and a finite nonnegative
measure $\zeta$ on $S$, such that
\begin{equation}\label{eq:variation:multiplier}
 2k\omega\mathcal E'(h,p)-c\int_{\mathfrak T}a'_C(h)\,d\pi(C)
 =\int_{\mathfrak A}\mathfrak C'(h,p;v)\,d\Lambda(x,v)
       -\int_S\theta'_+(h,p)\,d\zeta.
\end{equation}
This identity holds for all smooth compactly supported variations,
including arbitrary jets at $S$, and for the energy and momentum
prototypes specified below.  The scalar and vector moments of $\Lambda$
obey, initially as measures,
\begin{equation}\label{eq:variation:causality}
 u\ge |X|_g,\qquad u\mu+J(X)=0.
\end{equation}
Here a field and its density relative to the fixed original volume
form are identified when that density exists.
\end{proposition}

\begin{proof}
Choose $0<\beta<\min(1,q)$, let $\varrho$ be a smooth positive
extension of the end radius, and set $\rho_0=\varrho^{-n-\beta}$.  The strict conformal direction in Lemma~\ref{lem:strict-input} supplies
$\phi=O_6(r^{2-n})$, with finite flux,
\begin{equation}\label{eq:variation:strict-direction}
 -\Delta\phi=B\sqrt{|d\phi|^2+\varrho^{-2k}}+\rho_0,\qquad
 \partial_\nu\phi=-1,\qquad
 \frac{-\Delta\phi}{\rho_0}\longrightarrow1,
\end{equation}
where $B\ge|K|$ is smooth and $B=O_5(r^{-1-q})$.
The strict direction is $(2\phi g,\phi K)$.  The conformal variation
formula gives
\[
 \mathfrak C'=-4\phi(\mu+J(v))
       +2k\{-\Delta\phi+K(\nabla\phi,v)\}.
\]
It is strictly positive on active rays, its quotient by $\rho_0$ tends
to $2k$, and $-\theta'_+=k$ on $S$.

Take the real vector space generated by all compact variations, this
strict direction, and cut-off end prototypes of the following forms:
\begin{equation}\label{eq:variation:prototypes}
 h=r^{2-n}\delta,\ p=0;\qquad
 h=0,\quad p-(\tr_\delta p)\delta=\Pi(a),\quad
 \Pi(a)_{ij}=r^{-n}(a_ix_j+a_jx_i-(a\cdot x)\delta_{ij}).
\end{equation}
The flat scalar linearization of the first is zero outside the cutoff,
and $\partial^j\Pi(a)_{ij}=0$.  Their ADM derivatives are respectively
$E'=(n-2)/2$ and $P'_i=a_i/k$; thus they span all charge derivatives.
In the original background their constraint derivatives, also with
ray transport, are $O(r^{-n-q})$.  Their normalized values on the
active rays therefore tend to zero.

The coefficient of the strict direction is well defined.  Otherwise,
outside a compact set we would have $2a\phi g=b r^{2-n}\delta$ for
constants $a\ne0$ and $b$.  Taking the Euclidean trace gives
\[
 \phi=\frac{bn}{2a}\frac{r^{2-n}}{\tr_\delta g}
      =\frac b{2a}r^{2-n}+O_6(r^{2-n-q}).
\]
Since $\Delta_\delta r^{2-n}=0$, this would imply
$\Delta_g\phi=O(r^{-n-q})$, contradicting
$(-\Delta_g\phi)/\rho_0\to1$ and $\beta<q$.
This argument also covers bounded or empty active-ray sets.

Adjoin one infinity point to $\mathfrak A$, using its one-point
compactification when it is unbounded and an isolated point otherwise.
On this point assign $2k$ times the strict-direction coefficient.  The
linear map
\begin{equation}\label{eq:variation:separation-map}
 (h,p)\longmapsto
 \left(\frac{\mathfrak C'}{\rho_0},\ -\theta'_+,\
       \{c a'_C(h)-2k\omega\mathcal E'(h,p)\}_{C\in\mathfrak T}\right)
\end{equation}
takes values in continuous functions on the disjoint compact union
$\mathfrak A^\infty\sqcup S\sqcup\mathfrak T$.  Continuity on the last
factor is the strict tangent-plane continuity proved above.

Its image contains no strictly positive function.  To check this
without a constraint-surjectivity assumption, suppose such a variation
has strict-direction coefficient $a>0$, and write its remaining part
as $(h_0,p_0)$.  Use the path
\begin{equation}\label{eq:variation:feasible-path}
 g_t=e^{2at\phi}(g+th_0),\qquad
 K_t=e^{at\phi}(K+tp_0),\qquad t\ge0.
\end{equation}
On a compact region, strict positivity on the compact active set,
continuity on the unit-ball bundle, and Taylor's formula make all
constraints nonnegative for sufficiently small $t$.  This argument
also covers nearby inactive rays; the remaining compact inactive set
has a positive original minimum.

On the end the conformal identity leaves a positive multiple of the
nonnegative original contraction unaltered in sign.  The additional
non-strict linear error is $O(tr^{-n-q})$.  Quadratic prototype and
conformal errors are $O(t^2r^{2-2n})=O(t^2\rho_0)$, since $\beta<1$.
The strict conformal term is $2ka t\rho_0(1+o(1))$, uniformly in the
transported unit ball.  Choose the radius first, then $t$; it dominates
both errors and proves DEC there.  Likewise $\theta_+(S)<0$ for small
positive $t$.  Uniform metric equivalence preserves completeness and
positive full-cut area, and the perturbations preserve the stated
decay, integrability, and differentiable ADM limits.  Future
timelikeness of the charges is open.  Thus the numerical theorem applies
to Equation~\eqref{eq:variation:feasible-path}.

The third component of Equation~\eqref{eq:variation:separation-map}, together
with Equation~\eqref{eq:variation:envelope-derivative}, would give
$2k\omega\mathcal E'<c a'(0+)$, contradicting
$\mathcal E(g_t,K_t)\ge F(a(g_t))$ and equality at $t=0$.
This proves disjointness from the open positive cone.  Hahn--Banach
separation of an open convex cone from a linear subspace gives a
nonzero continuous functional, nonnegative on nonnegative functions,
annihilating the image.  No closed range is required.  By the Riesz
representation theorem it consists of three finite nonnegative
measures.

The mass of the measure on $\mathfrak T$ is positive.  Otherwise its
value on the strict direction would be positive, because the first
two components of that direction have positive minima and some
nonzero measure would remain on them.  Normalize this mass to one,
calling that measure $\pi$.  On the finite part of $\mathfrak A^\infty$
divide its measure by $\rho_0$, giving $\Lambda$, and call the boundary
measure $\zeta$.  An infinity atom contributes zero on compact
variations and prototypes, exactly the tests asserted in
Equation~\eqref{eq:variation:multiplier}.  It is therefore absent from that
identity, without having been assumed to vanish on the strict test.
The normalized identity is Equation~\eqref{eq:variation:multiplier}.

Define the moments by
$u(f)=\int f(x)\,d\Lambda$ and
$X(\alpha)=\int\alpha_x(v)\,d\Lambda$.
The unit-ball condition gives $|X|\le u$, and support on $\mathfrak A$
gives the complementarity identity in
Equation~\eqref{eq:variation:causality}.
\end{proof}

\subsection{Exact adjoints and regularity of the multipliers}

The separating measures are initially only measures. Their exact
adjoint equations supply the regularity and the positive Hessian
source needed to analyze the frontiers carrying area mass. The
modified-adjoint and null-perfect-fluid framework of
\cite[Section~6.1]{HuangLee2024Bartnik} is a useful precedent, but
does not supply these area-source terms or their support analysis.

\begin{lemma}[Interior equations]\label{lem:variation:adjoints}
In the open original exterior the moments have representatives
$u\in C^{0,1}_{\rm loc}$ and $X\in C^{1,1}_{\rm loc}$, and $u$ is
locally semiconvex.  Define the tensor-valued measures
\[
 d\sigma=\int_{\mathfrak T}dA_C\,d\pi(C),\qquad
 d\mathcal N=\int_{\mathfrak T}\nu_C^\flat\otimes\nu_C^\flat\,dA_C\,d\pi(C),
 \qquad d\mathcal T=g\,d\sigma-d\mathcal N.
\]
Their interior equations are
\begin{align}
 \sym\nabla X&=-uK,\label{eq:variation:shift}\\
 \Hess u-(\Delta u)g+\mathcal S(u,X,\nabla X)
      &=-\frac c2\,d\mathcal T,\label{eq:variation:metric-adjoint}\\
 \Hess u&=\mathcal B(u,X,\nabla X)+\frac c2\,d\mathcal N,
 \qquad
 \mathcal B=-\mathcal S+\frac{\tr\mathcal S}{k}g,
 \label{eq:variation:hessian}
\end{align}
where the smooth-coefficient linear expression is
\begin{equation}\label{eq:variation:S}
 \begin{split}
 \mathcal S={}&-u\Ric+2u(K^2-\tau K)
       +\mathcal L_XK-(\divg X)K\\
     &-\divg(K(X,\cdot))g-J_{(i}X_{j)}.
 \end{split}
\end{equation}
In particular the measure in Equation~\eqref{eq:variation:hessian} is positive
semidefinite.  Equations~\eqref{eq:variation:shift} and
\eqref{eq:variation:hessian} close after one prolongation on the first
jets of $u,X$, apart from this explicitly displayed measure.
\end{lemma}

\begin{proof}
All pairings in this proof use the fixed original volume density.
Direct differentiation gives
\begin{align}
 (2\mu)'={}&\nabla^i\nabla^jh_{ij}-\Delta\tr h-\langle\Ric,h\rangle
    +2\tau(\tr p-\langle K,h\rangle)
    -2\langle K,p\rangle+2\langle K^2,h\rangle,
 \label{eq:variation:scalar-linearization}\\
 J'_i={}&\nabla^jp_{ij}-\nabla_i\tr p
   +h^{jl}(\nabla_iK_{jl}-\nabla_lK_{ij})
   +\tfrac12K^{jl}\nabla_i h_{jl}
   -K_i{}^a\nabla^j h_{aj}
   +\tfrac12K_i{}^a\nabla_a\tr h,
 \label{eq:variation:momentum-linearization}\\
 \mathfrak C'={}&(2\mu)'+2J'(v)-J_i h^i{}_jv^j.
 \label{eq:variation:ray-linearization}
\end{align}
For Equation~\eqref{eq:variation:momentum-linearization}, vary both inverse
metrics and both connection terms in $\divg(K-\tau g)$; terms involving
$h\,d\tau$ and $\tau\,dh$ cancel.  The last term of
Equation~\eqref{eq:variation:ray-linearization} is the metric-square-root ray
transport, and must be retained.

On compact interior tests the objective derivative is zero.  Integration
by parts of the $p$ terms in Equation~\eqref{eq:variation:multiplier} gives
\[
 2\{u(\tau g-K)-\sym\nabla X+(\divg X)g\}=0.
\]
Taking its trace yields $\divg X=-u\tau$, and substitution proves
Equation~\eqref{eq:variation:shift}.  The curvature principal part has adjoint
$\Hess u-(\Delta u)g$.  Integrating the metric terms of $2J'(X)$ gives
\[
 \mathcal L_XK-(\divg X)K-\divg(K(X,\cdot))g.
\]
The algebraic scalar terms and ray transport give the other terms of
Equation~\eqref{eq:variation:S}.  The area derivative is
$\tfrac12\langle d\mathcal T,h\rangle$, proving
Equation~\eqref{eq:variation:metric-adjoint}.  Its trace is
$-k\Delta u+\tr\mathcal S=-ck\,d\sigma/2$;
substitution proves Equation~\eqref{eq:variation:hessian}, including its sign and
coefficient.

For the regularity assertion, differentiate
Equation~\eqref{eq:variation:shift} and commute derivatives.  One convenient
formula, in which the commutator convention is explicit, is
\begin{equation}\label{eq:variation:prolongation}
 \begin{split}
 \nabla_i\nabla_jX_l={}&-\nabla_i(uK_{jl})-\nabla_j(uK_{il})
                         +\nabla_l(uK_{ij})\\
 &+\tfrac12\big([\nabla_i,\nabla_j]X_l
       -[\nabla_i,\nabla_l]X_j-[\nabla_j,\nabla_l]X_i\big).
 \end{split}
\end{equation}
The commutators are smooth curvature coefficients times $X$.  Taking
traces of Equation~\eqref{eq:variation:hessian} and
Equation~\eqref{eq:variation:prolongation} gives a system with diagonal
Laplace principal part and smooth first-order lower terms.  Its sole
measure source has growth $O(r^{n-1})$, uniformly on compact subsets,
by the perimeter bound in Lemma~\ref{lem:variation:envelope}.

To spell out the regularity furnished by this growth, a local
order-minus-two parametrix for this smooth-coefficient elliptic system
has kernel bounded by $C|x-y|^{2-n}$ and first derivative bounded by
$C|x-y|^{1-n}$; its localized remainder is smooth.  These estimates
follow by freezing the principal matrix in a small coordinate ball,
using the Newton kernel, and successively correcting the smooth
coefficient error; first-order terms have one additional integrable
power.  Applied to a measure with mass at most $Cr^{n-1}$, the near part
of a first difference is $O(\delta)$ and the annuli between $\delta$
and a fixed radius contribute $O(\delta\log(1/\delta))$.  Hence the
potential is $C^{0,\alpha}$ for every $\alpha<1$.  The same local
parametrix applies to distributional solutions, leaving only a smooth
homogeneous remainder.  Thus both $u$ and $X$ are $C^{0,\alpha}$.
Equation~\eqref{eq:variation:prolongation}, or its trace in divergence
form, and the scalar Schauder estimate with Hölder divergence data
then give $X\in C^{1,\alpha}$.

Now $\mathcal B$ is locally bounded and continuous.  The positive tensor
in Equation~\eqref{eq:variation:hessian} implies a lower bound on the covariant
Hessian of $u$.  One can pass from this distributional statement to
semiconvexity without already assuming a bounded gradient.  In a smooth
chart mollify $u$.  The commutator between $\Gamma\,du$ and convolution
is bounded by $C\varepsilon^\alpha$: integrate its derivative onto the
kernel, subtract the value of $u$ at the center, and use the smooth
coefficient oscillation and the $C^{0,\alpha}$ bound.  Thus the
mollified covariant Hessian is bounded below uniformly.  Restriction
to a unit-speed geodesic gives a uniform lower bound for the ordinary
second derivative.  Uniform convergence transfers this bound to $u$.
Bounds on $u$ on slightly longer geodesic segments bound both one-sided
slopes on shorter segments; hence $u$ is locally Lipschitz.  Equation~\eqref{eq:variation:prolongation} now has a bounded right side in every
component, so $\nabla^2X\in L^\infty$ and $X\in C^{1,1}$.  Finally
bounded $du$ converts the covariant Hessian lower bound to ordinary
coordinate semiconvexity.  This also proves the stated closure on
first jets.
\end{proof}

\subsection{Compact normal variations through the active constraints}

The multiplier identity still contains every minimizing frontier. To
understand its typical support we need variations that move the metric
normally, $h=2sK$, while making the derivative of each active constraint
vanish. A Gaussian collar supplies such infinitesimal variations. It
is only a device for calculating these tests; no spacetime development
or energy condition for neighboring slices is assumed.


\begin{lemma}[Compact dust tests]\label{lem:variation:dust-tests}
Let $g,K$ be smooth on an open set $U$ of spatial dimension $n\geq3$, and
suppose that $\mu\geq |J|_g$, with the constraint and future-normal
conventions of this paper.  If $s\in C_c^\infty(U)$, there are smooth,
compactly supported tensor variations $(h,p_\delta)$, $\delta>0$, such that
\[
 h=2sK,\qquad
 \sup_{\substack{x\in U,\ |v|_g\leq1\\ \mu(x)+J_x(v)=0}}
 |\mathfrak C'(h,p_\delta;v)|\longrightarrow0.
\]
The supremum may equivalently be taken over a fixed compact neighborhood
of $\operatorname{supp}s$, since the variations vanish outside that
neighborhood.  Here the ray is transported with derivative
$v'=-\tfrac12h^\sharp v$, as in Equation~\eqref{eq:variation:multiplier}.
In particular, if the compact multiplier identity
Equation~\eqref{eq:variation:multiplier} holds with a locally finite active-ray
measure $\Lambda$, then these tests have
$\int\mathfrak C'(h,p_\delta;v)\,d\Lambda\to0$.
\end{lemma}

\begin{proof}
Fix relatively compact open neighborhoods of $\operatorname{supp}s$
whose closures lie in $U$.  All uniform limits below are on the larger
of these fixed neighborhoods.  Put
\[
 D_\delta=\frac{J\otimes J}{\mu+\delta}.
\]
We first realize these tensors as the spatial Einstein components of a
smooth Gaussian collar.  Write
\[
 \mathbf g_\delta=-dt^2+g_\delta(t),\qquad
 g_\delta(t)=g+2tK+t^2Q_\delta.
\]
For each fixed $\delta$ this is Lorentzian on a sufficiently short
collar.  Its future unit normal is $\mathbf n=\partial_t$, and the
initial second form is $K$.  At $t=0$, direct contraction of the Gaussian
Christoffel symbols
\[
 \boldsymbol\Gamma^0_{ij}=K_{ij},\qquad
 \boldsymbol\Gamma^i_{0j}=K^i{}_j,\qquad
 \boldsymbol\Gamma^i_{jk}=\Gamma^i_{jk}(g)
\]
gives
\[
 \begin{split}
 \mathbf{Ric}_{00}&=-\operatorname{tr}_gQ_\delta+|K|^2,\\
 \mathbf{Ric}_{0i}&=J_i,\\
 \mathbf{Ric}_{ij}&=\operatorname{Ric}_{ij}(g)
              +(Q_\delta)_{ij}+\tau K_{ij}-2(K^2)_{ij},\\
 \mathbf R&=R_g+2\operatorname{tr}_gQ_\delta
                   +\tau^2-3|K|^2.
 \end{split}
\]
Consequently $\mathbf G(\mathbf n,\mathbf n)=\mu$,
$\mathbf G(\mathbf n,\partial_i)=J_i$, and
\[
 \mathbf G_{ij}=(Q_\delta)_{ij}
                  -(\operatorname{tr}_gQ_\delta)g_{ij}+T_{ij},
 \quad
 T=\operatorname{Ric}_g+\tau K-2K^2
       -\tfrac12(R_g+\tau^2-3|K|^2)g.
\]
The map $Q\mapsto Q-(\operatorname{tr}_gQ)g$ is invertible, with
inverse $H\mapsto H-(\operatorname{tr}_gH)g/(n-1)$.  Thus the explicit
choice
\[
 Q_\delta=D_\delta-T
             -\frac{\operatorname{tr}_g(D_\delta-T)}{n-1}g
\]
has $\mathbf G_{ij}=D_\delta$.  This construction fixes every first
spacetime metric jet independently of $\delta$.  Collar widths and
higher normal derivatives need not have uniform bounds.

The tensors $D_\delta$ converge in spatial $C^1$ to
\[
 D_0=\begin{cases}J\otimes J/\mu,&\mu>0,\\0,&\mu=0.
 \end{cases}
\]
Here is the verification at the zero set, where division alone would
not justify the assertion.  Smoothness and nonnegativity give
$d\mu=0$ on $\{\mu=0\}$.  In any smooth local frame,
$|J_i|\leq C\mu$ implies $dJ_i=0$ there as well.  Differentiating the
displayed formula gives, with tensor norms taken using $g$,
\[
 |D_\delta|\leq\mu,\qquad
 |\nabla D_\delta|
 \leq 2|\nabla J|+|\nabla\mu|.
\]
On a small neighborhood of the compact zero set the right sides are
uniformly small.  On its complement $\mu$ has a positive lower bound,
so the formulas and their first derivatives converge uniformly.
Moreover $|D_0|\leq\mu$ makes $D_0$ differentiable, with zero derivative,
at each zero of $\mu$; the derivative bound makes this derivative
continuous.  These facts prove the asserted $C^1$ convergence.

Let $\mathbf T_\delta$ denote the Einstein tensor of the collar, and
identify its restriction to $t=0$ with a bilinear form on
$\mathbb R\mathbf n\oplus TU$.  Its time-zero components converge
spatially in $C^1$ to a tensor $\mathbf T_0$.
At a point where $\mu>0$ it is
\[
 \mathbf T_0=\mu^{-1}\alpha\otimes\alpha,
 \qquad \alpha(\mathbf n)=\mu,\quad \alpha|_{TU}=J.
\]
The covector $\alpha$ is causal in the precise sense that
$\alpha(a\mathbf n+w)\geq0$ when $a\geq|w|_g$ and $a\geq0$.
At an active ray with $\mu>0$, the inequalities
\[
 0=\mu+J(v)\geq\mu-|J|\,|v|\geq0
\]
force $|J|=\mu$, $|v|=1$, and $J=-\mu v^\flat$.
Thus $\ell=\mathbf n+v$ is null and $\alpha(\ell)=0$.
Parallel transport $\ell$ along any smooth curve in $\{t=0\}$ through
the point, using the collar connection, whose first jets are fixed.
The transported vector remains future null.  The smooth scalar
$\alpha(\ell)$ is nonnegative and vanishes at the point, so its
derivative in every spatial direction vanishes there.  Differentiating
$\mu^{-1}\alpha\otimes\alpha$ therefore gives
\[
 (\boldsymbol\nabla_Y\mathbf T_0)(W,\ell)=0
 \quad (Y\in TU,\ W\in\mathbb R\mathbf n\oplus TU).
\]
At $\mu=0$, all the components of $\mathbf T_0$ and all their spatial
first derivatives vanish.  Hence the same identity holds for every
active $|v|\leq1$ there, including the timelike vectors $\mathbf n+v$.

Contracted Bianchi supplies exactly the normal derivative which is
needed, without requiring convergence of third normal metric jets.
For a spatial orthonormal frame $e_1,\ldots,e_n$, evaluated at $t=0$,
\[
 (\boldsymbol\nabla_{\mathbf n}\mathbf T_\delta)
                  (\mathbf n,\ell)
   =\sum_{i=1}^n(\boldsymbol\nabla_{e_i}\mathbf T_\delta)(e_i,\ell).
\]
The right side converges uniformly to zero on the compact active-ray
set: spatial $C^1$ convergence was proved above, and all connection
coefficients in this formula are fixed at $t=0$.

Finally use the graph embeddings $F_a(x)=(a s(x),x)$ in this collar.
Differentiating the induced metric and future second form gives
\[
 h=2sK,\qquad p_\delta=\nabla^2s+sQ_\delta.
\]
Both variations are smooth and compactly supported.  We spell out the
argument-vector variation in order to include the terms involving
$ds$.  Transport $v$ on the domain by $v'_0=-sK^\sharp v$, and write
$\ell_a=\mathbf n_a+(F_a)_*v_a$.  Covariant differentiation along the
graph variation gives
\[
 \dot{\mathbf n}=\nabla s,\qquad
 \frac{D}{da}(F_a)_*v_a\bigg|_{a=0}=v(s)\mathbf n,
 \qquad \dot\ell=\nabla s+v(s)\mathbf n.
\]
In the middle equality, the spatial connection contribution
$sK^\sharp v$ cancels the prescribed ray transport.
Gauss--Codazzi on each graph identifies the differentiated ray
constraint with
\[
 \begin{split}
 \tfrac12\mathfrak C'(h,p_\delta;v)
  ={}&s(\boldsymbol\nabla_{\mathbf n}\mathbf T_\delta)
                      (\mathbf n,\ell)
       +\mathbf T_\delta(\nabla s,\ell)\\
     &+\mathbf T_\delta(\mathbf n,\nabla s+v(s)\mathbf n).
 \end{split}
\]
The first term tends uniformly to zero by Bianchi.  The second tends
to zero because $\mathbf T_0(\cdot,\ell)=0$ at every active ray.
The final term is exactly $J(\nabla s)+\mu v(s)$; it is zero at a
positive-density active ray by $J=-\mu v^\flat$, and at a zero-density
ray by $J=0$.  This proves the uniform assertion.  Local finiteness of
$\Lambda$ proves the final statement of the lemma.  Only infinitesimal
compact tests have been constructed; no energy condition on the
neighboring graphs is needed for this use of the multiplier identity.
\end{proof}


\subsection{Localizing a smooth family of minimizing enclosures}

Compact normal tests give an averaged trace identity. To turn it into
information about individual minimizing frontiers, we use their common
tangent plane at intersections. Positivity then prevents cancellation
between different frontiers. The remaining issue is contact with the
original boundary: the zero-expansion equation continues through contact,
and strong comparison identifies the boundary components that remain.

The next proposition takes the enclosure geometry as an explicit input.
Its regularity assumptions hold for the minimizing families considered
here in dimensions at most seven. They are not consequences merely of
smoothness of the original boundary in higher dimensions.

\begin{proposition}[Localization for smooth minimizing frontiers]
\label{prop:variation:smooth-localization}
Let $(\Omega,g,K)$ be smooth initial data on a connected manifold with
one end and nonempty compact smooth boundary
$S=\bigsqcup_{i=1}^r S_i$, where $r<\infty$ and each $S_i$ is connected.
Orient $S$ into $\Omega$, and suppose
\[
 H_S+\operatorname{tr}_S K=0,\qquad A=\operatorname{Area}_g(S)>0.
\]
A full enclosing cut means the entire compact boundary of a connected
closed exterior domain containing the sufficiently distant end, with
its manifold interior in $\operatorname{int}\Omega$. Its normal points
toward that end, and all components and all coincidence with $S$ count.
Assume $S$ is outer area-minimizing for these smooth cuts.

Let $\mathcal T$ be a measurable family of full minimizing frontiers,
equipped with a probability measure $\varpi$. Assume the following
geometric properties.
\begin{enumerate}
\item Every $T\in\mathcal T$ has area $A$, lies in one fixed compact
subset of $\Omega$, and is a compact embedded two-sided
$C^{1,\alpha}$ hypersurface without boundary, for some $0<\alpha<1$.
It is the entire boundary of a connected closed $C^1$ exterior domain
containing the distant end. Its manifold interior lies in
$\operatorname{int}\Omega$, and the coorientation of $T$ points toward
that end. In particular, whenever $T$ is smooth it is a full enclosing
cut of the stated comparison class.
\item The free part $T\setminus S$ is smooth and minimal. At contact
with $S$, its graph functions are locally $W^{2,2}$, its filled side
contains the original obstacle, and its normal agrees with that of $S$.
Thus its local graph lies on the exterior side of the obstacle graph.
\item The surface-area measures form a measurable kernel in $T$.
On the union of the free frontiers there is a single Borel tangent-plane
field $\Pi(x)$, continuous relative to that union, such that
$\Pi(x)=T_xT$ whenever $x\in T\setminus S$.
\end{enumerate}
Suppose also that
\begin{equation}
 \int_{\mathcal T}\int_T s\,\operatorname{tr}_{T_xT}K(x)
                   \,dA_g(x)\,d\varpi(T)=0
 \qquad\bigl(s\in C_c^\infty(\operatorname{int}\Omega)\bigr).
 \label{eq:variation:smooth-aggregate-trace}
\end{equation}
Then, for $\varpi$-almost every $T$, the whole frontier is a smooth
MOTS, its free part satisfies $H_T=\operatorname{tr}_T K=0$, and each
of its connected components is either an entire $S_i$ or is compactly
contained in $\operatorname{int}\Omega$.

Moreover, $\varpi\{T:T=S\}=1$ under either of the following hypotheses:
\begin{enumerate}
\item[(a)] $S$ is connected, and no compact smooth embedded two-sided
full enclosing cut lying wholly in $\operatorname{int}\Omega$, possibly
disconnected, has $\theta_+\le0$ everywhere.
\item[(b)] No full smooth enclosing cut has all components either equal
to entire original components $S_i$ or wholly in
$\operatorname{int}\Omega$, at least one component in that open exterior,
and $\theta_+\le0$ everywhere.
\end{enumerate}
In case \textup{(a)}, before outermostness is applied, almost every
frontier satisfies the sharper dichotomy: it equals $S$, or it is
wholly interior and satisfies $H_T=\operatorname{tr}_T K=0$ everywhere.
In either localization case, for every continuous symmetric tensor $h$,
\begin{equation}
 \int_{\mathcal T}\frac12\int_T\operatorname{tr}_{T_xT}h\,dA_g
                                  \,d\varpi(T)
       =\frac12\int_S\operatorname{tr}_S h\,dA_g.
 \label{eq:variation:smooth-localized-area}
\end{equation}
\end{proposition}

\begin{proof}
Define the aggregate area measure in the open exterior by
\[
 \beta(B)=\int_{\mathcal T}\operatorname{Area}_g(T\cap B)\,d\varpi(T)
 \qquad\bigl(B\subset\operatorname{int}\Omega\text{ Borel}\bigr).
\]
It is positive and has mass at most $A$. On the union of the free
frontiers put $F(x)=\operatorname{tr}_{\Pi(x)}K(x)$, and set $F=0$
elsewhere. This is a bounded measurable function on the common compact
region. Equation~\eqref{eq:variation:smooth-aggregate-trace} says that
the signed Radon measure $F\beta$ vanishes. Hence $F=0$
$\beta$-almost everywhere, and Tonelli's theorem gives
\[
 0=\int |F|\,d\beta
  =\int_{\mathcal T}\int_{T\setminus S}
          |\operatorname{tr}_{T_xT}K|\,dA_g\,d\varpi(T).
\]
For almost every $T$, its tangential trace therefore vanishes at
area-almost every free point. Each free part is smooth, so a nonzero
value would persist on an open patch of positive area. Thus its trace
vanishes everywhere there. Minimality gives $H_T=0$ on that part as
well. The common-plane hypothesis is essential: it makes the trace one
function of position before positivity is used.

Fix one of these frontiers and a contact point. Extend the smooth
one-sided coefficients of $g,K$ locally across $S$ only to write graph
equations. In a graph chart let $f$ describe $T$ and let the smooth
function $\psi$ describe the obstacle. Their ordered graphs have
matching normals at contact. On the coincidence set $Z=\{f=\psi\}$,
their first derivatives agree. Sobolev locality applied to
$\partial_j(f-\psi)\in W^{1,2}$ gives
$D^2f=D^2\psi$ almost everywhere on $Z$. Since $S$ is marginal, the
graph of $f$ satisfies $H+\operatorname{tr}_{\rm tan}K=0$ almost
everywhere on $Z$. It satisfies this equation on the free set by the
first part of the proof.

We verify that this almost-everywhere equation yields smoothness
through contact. In the chart it has the form
\[
 \partial_i\mathcal F^i(y,f,Df)=\mathcal B(y,f,Df),
\]
with smooth functions $\mathcal F,\mathcal B$. The matrix
$\mathcal F^i_{p_j}$ is uniformly elliptic on a smaller chart, since
$Df$ is bounded there. These are the usual graph coefficients for
mean curvature, with the smooth tangential trace of $K$ included in
$\mathcal B$. As $f\in W^{2,2}$, the divergence on the left is an
$L^2$ function. The almost-everywhere identity is therefore also a
weak divergence-form equation; there is no additional contact measure.
For $f_k=\partial_k f$, differentiation in distributions gives
\[
 \partial_i(a^{ij}\partial_j f_k)
   =\partial_i\bigl(\delta_{ik}\mathcal B
           -\mathcal F^i_{y_k}-\mathcal F^i_z f_k\bigr),
 \qquad a^{ij}=\mathcal F^i_{p_j}(y,f,Df).
\]
All displayed coefficients and the right-hand vector field are
$C^{0,\alpha}$, while $f_k\in W^{1,2}$. Interior divergence-form
Schauder regularity gives $f_k\in C^{1,\alpha}$ on smaller charts.
Thus $f\in C^{2,\alpha}$, and further differentiation and Schauder
estimates give smoothness; see~\cite{GilbargTrudinger2001}.
Consequently $T$ is a smooth MOTS across every contact point.

If $T$ touches $S_i$, the two ordered smooth graphs satisfy the same
zero-expansion equation with the same orientation. Their nonnegative
difference $w$ satisfies a linear uniformly elliptic equation
$a^{ij}w_{ij}+b^i w_i+cw=0$ with bounded coefficients. Replacing $c$
by $\min(c,0)$ gives an inequality with left side $\le0$, since
$w\ge0$. The strong minimum principle at a zero of $w$ forces local
coincidence. Therefore $T\cap S_i$ is both open and closed in the
connected hypersurface $S_i$, so $S_i\subset T$. Local coincidence
also makes $S_i$ open in $T$, and compactness makes it closed there;
it is one connected component of $T$. Every remaining component is
disjoint from $S$ and, by compactness, lies a positive distance inside
the open exterior. This proves the component assertion.

Suppose first that \textup{(a)} holds. If $T$ touches $S$, it contains
$S$ as a component. Since $\operatorname{Area}_g(T)=A=
\operatorname{Area}_g(S)$, no further component of positive area is
possible, and $T=S$. Otherwise $T$ is a wholly interior smooth full
enclosing MOTS, excluded by \textup{(a)}.

Under \textup{(b)}, any interior component of $T$ would give exactly
a prohibited partial-coincidence cut: all its other components are
interior or entire $S_i$, and $\theta_+=0$ everywhere. Hence
$T=\bigsqcup_{i\in I}S_i$ for some $I\subset\{1,\ldots,r\}$.
The area identity
\[
 \sum_{i\in I}\operatorname{Area}_g(S_i)
       =A=\sum_{i=1}^r\operatorname{Area}_g(S_i)
\]
and the positive area of every $S_i$ force $I=\{1,\ldots,r\}$.
Thus $T=S$ in this case too. All conclusions hold outside the single
null family discarded above. Integrating the area derivative of this
fixed frontier proves Equation~\eqref{eq:variation:smooth-localized-area}.
\end{proof}


\subsection{The typical minimizing frontiers}

\begin{lemma}[Zero trace and contact dichotomy]
\label{lem:variation:contact-dichotomy}
For $\pi$-almost every $C\in\mathfrak T$, either $\Gamma_C=S$, or
$\Gamma_C$ is compactly contained in the open exterior and
\begin{equation}\label{eq:variation:three-equations}
 H_{\Gamma_C}=\tr_{T\Gamma_C}K=0
\end{equation}
on its regular part.  The latter frontier is not asserted to be smooth
at this stage.
\end{lemma}

\begin{proof}
Apply Lemma~\ref{lem:variation:dust-tests} to an arbitrary compact
smooth lapse $s$ in the open exterior.  The variations have zero ADM
and boundary terms.  Their active-constraint derivatives tend uniformly
to zero on a fixed compact set, where $\Lambda$ has finite mass.
Equation~\eqref{eq:variation:multiplier} therefore gives
\begin{equation}\label{eq:variation:trace-measure}
 \int_{\mathfrak T}\int_{\partial^*C}
           s\,\tr_{T\partial^*C}K\,dA\,d\pi(C)=0.
\end{equation}
On a uniform graph cylinder from
Lemma~\ref{lem:variation:envelope}, all sheets through one point have
the same tangent plane.  Their tangent planes also vary continuously
on the union of the sheets: a contrary sequence would contradict
multiplicity-one regular convergence.  Consequently
$\tr_{T\partial^*C}K$ defines one continuous function $f$ on this union,
and Equation~\eqref{eq:variation:trace-measure} says $f\,d\sigma=0$ there.
There is no cancellation between differing tangent planes at the same
point.  Hence $f=0$ almost everywhere for the positive measure
$\sigma$.  Tonelli's theorem and continuity on each smooth graph give
$\tr_{T\Gamma_C}K=0$ everywhere on each participating free graph for
almost every $C$.  Choose countably many smaller cylinders covering
all free regular points of the family and discard the union of their
null exceptional sets.  Minimality already gives $H=0$ on these
patches.

Suppose a remaining frontier touches $S$.  Near a contact point write
the frontier and $S$ as $W^{2,p}\cap C^{1,\alpha}$ and smooth ordered
graphs, respectively.  Their difference vanishes on the coincidence
set.  Sobolev locality says that its first and second weak derivatives
vanish almost everywhere on that set; this follows by applying
first-derivative locality to the function and then to its first
derivatives.  Thus the graph's second jet agrees there almost
everywhere with that of $S$.  On the coincidence set it has
$H+\tr_{\rm tan}K=\theta_+(S)=0$, and on its free set it has the same
equation by Equation~\eqref{eq:variation:three-equations}.  It therefore solves
the uniformly elliptic scalar zero-expansion graph equation almost
everywhere on the entire patch.  Linear $W^{2,p}$ estimates and
Schauder bootstrap make it smooth.  The strong comparison principle
for two ordered solutions of this graph equation makes it coincide
locally with $S$.  Connectedness of $S$, local continuation of this
coincidence, and closedness of the frontier make it contain all of
$S$.  Since its total perimeter is $A=\Area(S)$, there is no remaining
perimeter for another sheet.  The support is exactly $S$.

Otherwise its compact support is disjoint from the compact set $S$,
so it lies a positive distance inside the open exterior and satisfies
Equation~\eqref{eq:variation:three-equations} on every regular patch.
\end{proof}

When $3\leq n\leq7$, Proposition~\ref{prop:variation:smooth-localization}
now gives the desired localization directly. Indeed,
Lemma~\ref{lem:variation:envelope} supplies compact minimizing frontiers
of area $A$, their full enclosing geometry, smooth minimal free parts,
and the required contact regularity. Its common-plane property and
multiplicity-one regular convergence give the continuous plane field.
Equation~\eqref{eq:variation:trace-measure} is the averaged trace identity.
The connected marginal boundary and wholly-interior outermostness in
Definition~\ref{def:setting:rigidity} match case \textup{(a)} of the
proposition. Thus $\pi\{C:\Gamma_C=S\}=1$.
In these dimensions the reader may proceed to
Theorem~\ref{thm:variation:boundary-support} and then to
Section~\ref{sec:static}. The weaker-decay classes in dimensions three
and four require their own end analysis below; this application uses
the strong-decay class of Definition~\ref{def:data}.


For $n\geq8$, smooth outermostness cannot yet be applied: the wholly
interior frontier could have singular points. The next arguments prove
that almost every frontier carrying area mass is nevertheless smooth.
We distinguish atoms of the local graph-height measure from its
nonatomic part.
An atom gives a nonzero jump of the normal derivative of $u$; the
tangential Hessian equation then forces that sheet to be totally
geodesic. For a nonatomic sheet, neighboring minimizers produce one
common positive Jacobi field. We will show that this field diverges
at every singular end and that the lapse is positive somewhere on
almost every such sheet. Focusing then bounds its second fundamental
form, removes its singularities, and permits the original smooth
outermostness hypothesis to apply.

\subsection{Measurable graph families and the atomic case}
\label{subsec:variation:graph-families}

We first fix the measurable meaning of this local distinction. Choose
nested uniform graph cylinders
$Q_j^-\Subset Q_j^0\Subset Q_j^+$ such that every member meeting $Q_j^0$
has one graph over the fixed base used in $Q_j^-$, with uniform estimates
inside $Q_j^+$. Shrinking the cylinders from
Lemma~\ref{lem:variation:envelope}, and using their uniform slope bounds,
gives these margins. The smallest cylinders cover every free regular
point of the minimizing family; second countability gives a countable
subcover.

Put $\mathcal U_j=\{D\in\mathfrak T:\Gamma_D\cap Q_j^0\ne\varnothing\}$.
This is open in $\mathfrak T$: strict perimeter convergence and the
density lower bounds imply local convergence of perimeter supports.
It has a countable compact exhaustion, for example by the sets at
distance at least $1/m$ from $\mathfrak T\setminus\mathcal U_j$;
if $\mathcal U_j=\mathfrak T$, use $\mathfrak T$ itself.
Graph height at a fixed base point defines a continuous map
$a_j:\mathcal U_j\to\R$. Its image $I_j$ is therefore a countable union
of compact sets and is Borel. Identify graphs with equal height:
strong comparison makes them coincide, and
Equation~\eqref{eq:variation:graph-separation} makes the graph map
$a\mapsto h_a$ Lipschitz into each smaller $C^l$ graph space.
Thus the parameter family and its geometric weights are Borel.

Let $\lambda_j=(a_j)_*(\pi|_{\mathcal U_j})$ be its finite height
measure. Every sheet contributing in $Q_j^-$ is represented, including
a sheet crossing the edge of a larger cylinder. Distinct parameters
have distinct heights at every base point by strong comparison.
An \emph{atomic patch} means a graph whose parameter has positive
$\lambda_j$ mass; equivalently, the fiber of minimizers agreeing on
that patch has positive $\pi$ mass. The nonatomic alternative refers
to the nonatomic part of this same pushforward measure. All later
exceptional families will be discarded on the original space
$\mathfrak T$, using these countably many maps.

Figure~\ref{fig:variation:singular-localization} shows how the two
cases lead to the same regularity conclusion. For a connected component
$\Sigma$ of the regular frontier, write $A_\Sigma$ for its second
fundamental form and $K_{\rm tan}=K|_{T\Sigma}$ for the tangential
restriction of $K$. The nonatomic route
uses both the Jacobi field at singular ends and the diffuse Hessian
measure; neither input alone gives the required maximum.

\begin{figure}[htbp]
\centering
\begin{tikzpicture}[
 box/.style={draw,rounded corners=2pt,align=center,font=\small,
             inner sep=5pt,text width=.39\linewidth},
 flow/.style={-{Latex[length=2mm]},thick}]
 \node[box,text width=.83\linewidth] (family) at (0,0)
   {Wholly interior minimizing frontier, possibly singular\\
    $H=\tr_{\rm tan}K=0$ on its regular part};
 \node[box] (atom) at (-3.5,-1.35)
   {Atomic graph parameter\\
    normal derivative jump};
 \node[box] (nonatomic) at (3.5,-1.35)
   {Nonatomic graph parameter\\
    one common positive Jacobi field $\varphi$};
 \node[box] (flat) at (-3.5,-2.95)
   {Tangential Hessian traces\\
    $A_\Sigma=0$};
 \node[box] (maximum) at (3.5,-2.95)
   {Diffuse BV: $u>0$ somewhere\\
    Singular ends: $\varphi\to\infty$\\
    Maximum principle and focusing\\
    $A_\Sigma\pm K_{\rm tan}=0$ for one sign};
 \node[box,text width=.83\linewidth] (bound) at (0,-4.85)
   {Uniform curvature bound $|A_\Sigma|\le\sup_{\Gamma_C}|K|$\\
    on every regular component};
 \node[box,text width=.83\linewidth] (finish) at (0,-6.2)
   {Planar tangent cones remove singularities\\
    Smooth outermostness excludes the interior frontier};
 \draw[flow] (family.south) -- (atom.north);
 \draw[flow] (family.south) -- (nonatomic.north);
 \draw[flow] (atom) -- (flat);
 \draw[flow] (nonatomic) -- (maximum);
 \draw[flow] (flat.south) -- (bound.north);
 \draw[flow] (maximum.south) -- (bound.north);
 \draw[flow] (bound) -- (finish);
\end{tikzpicture}
\caption{Removal of singular frontiers carrying area mass. Each regular
component follows one of the two routes. The common curvature bound holds
for one full-measure family of minimizers, after the countable
exceptional-family step in
Theorem~\ref{thm:variation:boundary-support}.}
\label{fig:variation:singular-localization}
\end{figure}

\begin{lemma}[Atoms give totally geodesic sheets]\label{lem:variation:atomic-flatness}
Fix a uniform graph cylinder about a free regular sheet $\Sigma$ of a
minimizer.  Collapse coincident graphs and parametrize the remaining
graphs by their height at a fixed base point.  If the parameter of
$\Sigma$ has positive mass under the pushforward of $\pi$, then the
second fundamental form $A_\Sigma$ vanishes on its smaller patch.
\end{lemma}

\begin{proof}
Use Fermi coordinates $(y,s)$ about the smooth reference sheet, with
$g=ds^2+g_s$, $s=0$ on $\Sigma$, and
$A_{ab}=\tfrac12\partial_s(g_s)_{ab}$.  Let $a>0$ be the mass of the
coincident-graph fiber.  Its contribution to
Equation~\eqref{eq:variation:hessian} is
$(ca/2)\,ds\otimes ds\,dA_\Sigma$.
In the $ss$ equation this is the density $(ca/2)\delta_0(ds)$ on
almost every normal line; the volume and surface densities coincide
at $s=0$.  All other graph parameters have height different from zero
on every such line, by noncrossing and contact agreement.  They
contribute no atom at zero.  Semiconvexity and local Lipschitz bounds
give bounded one-sided traces of $\partial_su$ on almost every line,
and slicing the equation proves
\begin{equation}\label{eq:variation:normal-jump}
 (\partial_su)^+-(\partial_su)^-=ca/2.
\end{equation}

The tangential equations contain no jump measure on the reference
sheet.  Contributions from nearby sheets must also be controlled
before taking traces.  Equation~\eqref{eq:variation:graph-separation} implies that a nearby leaf at
distance $|s|$ has tangential normal components $O(|s|)$ relative to
the reference Fermi frame, on a smaller patch.  Indeed Harnack
compares its height at the fixed base point to its height at the
chosen point, and the derivative estimate controls its tilt.  Thus
the tangential--tangential components of $d\mathcal N$ in
$|s|<\varepsilon$ have total variation bounded by
$C\varepsilon^2\sigma(\{|s|<\varepsilon\})$.  After division by the
slab width this tends to zero; the total local $\sigma$ mass is
bounded.  This estimate also covers infinitely many accumulating
leaves and their diffuse parameter measure.

Average the tangential equation over $0<s<\varepsilon$ and over
$-\varepsilon<s<0$, and test in the $y$ variables.  The continuous
tensor $\mathcal B(u,X,\nabla X)$ has identical limiting traces.  On
each parallel sheet,
\[
 (\Hess_g u)_{ab}=(\Hess_{g_s}u|_{s})_{ab}
                              +A_{ab}(s)\partial_su.
\]
Uniform continuity of $u$ makes the two averaged intrinsic Hessians
converge to the same distribution, and the bounded normal derivatives
converge to their one-sided traces in $L^1$ on smaller base patches.
Taking the difference and using the vanishing measure contribution
therefore gives $A_{ab}(0)[\partial_su]=0$ as distributions.  Equation~\eqref{eq:variation:normal-jump} and smoothness of $A$ prove $A=0$.
\end{proof}

\subsection{Jacobi fields from neighboring minimizers}

On a free sheet with Equation~\eqref{eq:variation:three-equations}, let $L_0$
be the minimal Jacobi operator and $L_\pm$ the normal linearizations of
$H\pm\tr_{\rm tan}K$.  Explicitly,
\begin{align}
 L_0f&=-\Delta_\Sigma f-(|A_\Sigma|^2+\Ric(\nu,\nu))f,
 \label{eq:variation:L0}\\
 L_\pm f&=L_0f\ \pm2K(\nu,\nabla_\Sigma f)
                      \pm\tr_{T\Sigma}(\nabla_\nu K)f.
 \label{eq:variation:Lpm}
\end{align}
These formulas follow from $\nu'=-\nabla_\Sigma f$ and the
mean-curvature variation formula.  In the trace of $K$, the two terms
from differentiating the tangent frame cancel the variation of the
inverse induced metric, leaving precisely the last two terms above.

\begin{lemma}[One common positive Jacobi field]
\label{lem:variation:common-jacobi}
Let $C\in\supp\pi$ belong to the full-measure family in
Lemma~\ref{lem:variation:contact-dichotomy}, with wholly interior
frontier, and let $\Sigma$ be a connected component of its regular
part.  If its local graph parameter at some patch is not an atom,
there is a smooth positive function $\varphi$ on all of $\Sigma$ with
\begin{equation}\label{eq:variation:common-jacobi}
 L_0\varphi=L_+\varphi=L_-\varphi=0.
\end{equation}
\end{lemma}

\begin{proof}
Fix a base point $y_0$ in the patch.  The fiber of minimizers whose
graph agrees there has $\pi$ measure zero.  Every neighborhood of $C$
in the compact metric space $\mathfrak T$ has positive $\pi$ measure.
Deleting this fiber and the exceptional family from
Lemma~\ref{lem:variation:contact-dichotomy}, choose $C_j\to C$ with
distinct graphs at $y_0$, all satisfying
Equation~\eqref{eq:variation:three-equations} near the component.

Choose a connected relatively compact smooth exhaustion of $\Sigma$
by domains containing $y_0$.  For each fixed exhaustion domain,
multiplicity-one regular convergence and finitely many overlapping
graph charts express $\Gamma_{C_j}$ as one smooth normal graph $h_j$
over it, with $h_j\to0$ smoothly.  Noncrossing says this graph either
agrees with $\Sigma$ or is strictly on one side.  Agreement at one
point propagates through overlapping charts by strong comparison;
it would contradict distinction at $y_0$.  A subsequence thus has a
fixed strict sign on every exhaustion domain.

Divide $h_j$ by its signed value at $y_0$, obtaining positive functions
normalized to one there.  Each satisfies the linear difference
equation for the two minimal graphs.  The coefficients converge
smoothly to those of $L_0$.  Harnack along a finite chain of balls
and interior estimates bound the normalized graphs and all their
derivatives on each compact subset.  Diagonal extraction gives a
smooth positive limit $\varphi$ on all of $\Sigma$.  The same graph
differences solve the two difference equations for
$H\pm\tr_{\rm tan}K=0$; their coefficients converge to $L_\pm$.
Passing to the same limit proves all three equations in
Equation~\eqref{eq:variation:common-jacobi}.  This argument requires a sequence
of neighboring minimizers, not a differentiable family.
\end{proof}


\subsection{Focusing and the positive stress measure}

The common Jacobi field will provide the comparison function for a
maximum principle. We now establish the inequality it will be compared
with. The positive Hessian source in the lapse equation becomes a
positive spatial stress after trace reversal; retaining that source is
what gives the favorable sign in the focusing inequality.


\begin{lemma}[The stationary metric adjoint]
\label{lem:variation:stationary-action}
Let $u>0$ and $X$ be smooth on an open set with smooth metric $g$, and put
\[
 \mathbf g=-u^2dz^2+g_{ij}(dx^i+X^i dz)(dx^j+X^j dz),
 \qquad b=-u^{-1}\operatorname{sym}\nabla X.
\]
Let $\mu_b,J_b$ be the constraints of $(g,b)$, and let $S_b$ denote
the expression $\mathcal S$ in Equation~\eqref{eq:variation:S}, with
$K$ replaced throughout by $b$.  Then the exact smooth identity is
\[
 \nabla^2u-(\Delta u)g+S_b
  =-u\mathbf G_{\mathrm{spatial}}
        -(u\mu_b+J_b(X))g-(J_b)_{(i}X_{j)}.
\]
Here the spatial Einstein tensor is its restriction to the tangent
spaces of $z=\mathrm{constant}$, and parentheses include the factor
$1/2$.
\end{lemma}

\begin{proof}
All variations and integrations in this proof have compact support in
the open set.  Define
\[
 \mathsf B(B,C)=\langle B,C\rangle_g
                    -(\operatorname{tr}_gB)(\operatorname{tr}_gC).
\]
Temporarily let $K$ be independent of the metric.  Integration of the
momentum divergence, at fixed contravariant shift $X$, gives
\[
 \begin{split}
 \int(2u\mu+2J(X))\,dV_g
 &=\int u\{R_g-\mathsf B(K,K)+2\mathsf B(K,b)\}\,dV_g\\
 &=\int u\{R_g+\mathsf B(b,b)
                  -\mathsf B(K-b,K-b)\}\,dV_g.
 \end{split}
\]
At $K=b$ the final square has zero first variation, including for a
metric variation with the covariant components of $K$ held fixed.
Thus the metric derivative of the constraint functional at this point
equals that of the stationary scalar action
$\int u(R_g+\mathsf B(b,b))\,dV_g$.

For completeness, its relation to the spacetime scalar action follows
from the Gaussian contraction formulas with lapse and shift.  They
give
\[
 \mathbf R=R_g+|b|^2+(\operatorname{tr}b)^2
       +2\mathbf n(\operatorname{tr}b)-2u^{-1}\Delta u.
\]
Stationarity implies
$\mathbf n(\operatorname{tr}b)=-u^{-1}X(\operatorname{tr}b)$ and
$\operatorname{div}X=-u\operatorname{tr}b$, whence
\[
 \mathbf R=R_g+\mathsf B(b,b)
        -2u^{-1}\operatorname{div}
                    \bigl(\nabla u+(\operatorname{tr}b)X\bigr).
\]
Since $\sqrt{|\det\mathbf g|}=u\sqrt{\det g}$, the two scalar actions
differ by a spatial divergence per unit $z$-length.  For a spacetime
metric variation $H$, differentiating the connection formula for
scalar curvature gives
\[
 \delta(\mathbf R\,dV_{\mathbf g})
   =\left[-\langle\mathbf G,H\rangle_{\mathbf g}
    +\boldsymbol\nabla_\alpha
       \bigl(\boldsymbol\nabla_\beta H^{\alpha\beta}
              -\boldsymbol\nabla^\alpha\operatorname{tr}_{\mathbf g}H
       \bigr)\right]dV_{\mathbf g}.
\]
At fixed $u,X$, $H=h_{ij}(dx^i+X^i dz)(dx^j+X^j dz)$ annihilates the
unit normal in each slot.  The derivative of the stationary action
therefore pairs $h$ with $-u\mathbf G_{\mathrm{spatial}}$.

The metric coefficient of the constraint integral before ray
transport is
\[
 \nabla^2u-(\Delta u)g+S_b+(J_b)_{(i}X_{j)}
                   +(u\mu_b+J_b(X))g.
\]
The last term is the volume derivative, and the preceding added term
restores the term removed by isometric ray transport in the
definition of $S_b$.  Equating the two coefficients proves the
identity.  This also displays why the volume derivative disappears
only after complementarity has been imposed.
\end{proof}

\begin{lemma}[Focusing with a positive Hessian measure]
\label{lem:variation:measure-focusing}
Let $g,K$ be smooth on $U$, with $\mu\geq|J|_g$.  Suppose that $u$ is
locally Lipschitz and semiconvex, $X\in C^{1,1}_{\mathrm{loc}}$, and
the following exact distributional equations hold:
\[
 \begin{gathered}
 u\geq|X|_g,\qquad u\mu+J(X)=0,\qquad
 \operatorname{sym}\nabla X=-uK,\\
 \nabla^2u=-S+\frac{\operatorname{tr}_gS}{n-1}g+\mathcal M,
 \end{gathered}
\]
where
\[
 \begin{split}
 S={}&-u\operatorname{Ric}_g+2u(K^2-\tau K)
        +\mathcal L_XK-(\operatorname{div}X)K\\
     &-\operatorname{div}(K(X,\cdot))g-J_{(i}X_{j)},
 \end{split}
\]
and $\mathcal M$ is a locally finite positive semidefinite symmetric
tensor measure, interpreted relative to $dV_g$.  In the application,
\[
 \mathcal M=\frac c2\int\nu_D^\flat\otimes\nu_D^\flat\,dA_D\,d\pi.
\]
Let $\Sigma\subset U$ be a smooth two-sided hypersurface with
$H_\Sigma=\operatorname{tr}_\Sigma K=0$, and let $\nu$ be its chosen
unit normal.  Set
\[
 \begin{split}
 L_\pm f={}&-\Delta_\Sigma f
        \pm2K(\nu,\nabla_\Sigma f)\\
 &+\bigl[-|A_\Sigma|^2-\operatorname{Ric}_g(\nu,\nu)
                 \pm\operatorname{tr}_\Sigma(\nabla_\nu K)\bigr]f,
 \qquad Q_\pm=u\mp g(X,\nu).
 \end{split}
\]
Then $Q_\pm\geq0$, and, on the open subset $\Sigma\cap\{u>0\}$,
\[
 L_\pm Q_\pm\leq-u|A_\Sigma\pm K_{\mathrm{tan}}|^2
\]
in distributions.  Precisely, each inequality holds after pairing
with any nonnegative smooth compactly supported test function on
that open subset.
\end{lemma}

\begin{proof}
The operator formula follows by varying $\Sigma$ with velocity $f\nu$.
The mean-curvature derivative is
$-\Delta_\Sigma f-(|A_\Sigma|^2+\operatorname{Ric}_g(\nu,\nu))f$.
For the other term, differentiation of
$\operatorname{tr}_\Sigma K=\operatorname{tr}_gK-K(\nu,\nu)$, using
$\dot\nu=-\nabla_\Sigma f$, gives
$f\operatorname{tr}_\Sigma(\nabla_\nu K)
 +2K(\nu,\nabla_\Sigma f)$.  Adding the two derivatives gives $L_+$;
replacing $K$ by $-K$ gives $L_-$.

Fix a compact patch of $\Sigma\cap\{u>0\}$.  Enlarge it slightly to
a coordinate neighborhood on which $u\geq a>0$.  All constants below
may depend on this neighborhood, $a$, the smooth data, and the stated
local Lipschitz bounds; the only limiting parameter is the mollifier
radius $\varepsilon\downarrow0$.  Choose a nonnegative smooth
convolution kernel of integral one.  Mollify $u$ and the components
of the covector $w=X^\flat$, and then put
\[
 u_\varepsilon=u*\rho_\varepsilon,\qquad
 w_\varepsilon=w*\rho_\varepsilon,\qquad
 X_\varepsilon=g^{-1}w_\varepsilon,
 \qquad b_\varepsilon=-u_\varepsilon^{-1}
                    \operatorname{sym}\nabla X_\varepsilon.
\]
Convolutions are used only on smaller coordinate neighborhoods.
In particular $u_\varepsilon\geq a/2$ for small $\varepsilon$.

We need convergence of $b_\varepsilon$ through one derivative.
For a smooth coefficient $a_0$ and a Lipschitz function $f$, let
\[
 C_\varepsilon(a_0,f)=(a_0f)*\rho_\varepsilon
                              -a_0(f*\rho_\varepsilon).
\]
The integral formula with integrand
$(a_0(y)-a_0(x))f(y)\rho_\varepsilon(x-y)$ gives
$\|C_\varepsilon(a_0,f)\|_\infty\leq C\varepsilon\|f\|_\infty$.
Differentiation in distributions gives the useful identity
\[
 \partial_j C_\varepsilon(a_0,f)
    =C_\varepsilon(\partial_j a_0,f)
                         +C_\varepsilon(a_0,\partial_jf).
\]
The same integral bound, applied to the bounded weak derivative of
$f$, proves
\[
 \|C_\varepsilon(a_0,f)\|_{C^1}
       \leq C\varepsilon\|f\|_{W^{1,\infty}}.
\]
In covector coordinates the shift equation reads
$\tfrac12(\partial_iw_j+\partial_jw_i)-\Gamma^l_{ij}w_l=-uK_{ij}$.
Thus
\[
 \operatorname{sym}\nabla X_\varepsilon+u_\varepsilon K
   =C_\varepsilon(\Gamma,w)-C_\varepsilon(K,u)=O_{C^1}(\varepsilon).
\]
The positive lower bound for $u_\varepsilon$ and its bounded first
derivatives therefore give
\[
 b_\varepsilon\longrightarrow K\quad\hbox{in }C^1.
\]
Also $X_\varepsilon\to X$ in $C^1$, $u_\varepsilon\to u$ uniformly,
and all their first derivatives are uniformly bounded.

Put $B=-S+(\operatorname{tr}_gS)g/(n-1)$; it is continuous under the
stated regularity.  To convolve the tensor measure precisely, write
$dV_g=v_g\,dx$ in the chosen chart and convolve the component measures
$v_g^{-1}\,d\mathcal M_{ij}$.  Denote the resulting matrix of smooth
functions by $M_\varepsilon$.  This matrix is positive semidefinite at
each point, because the convolution kernel and $v_g^{-1}$ are
nonnegative.  The Hessian equation and the connection commutator give
\[
 \nabla^2u_\varepsilon=B*\rho_\varepsilon
                            +M_\varepsilon+o_{C^0}(1).
\]
Indeed the sole extra term is
$(\Gamma\,du)*\rho_\varepsilon-\Gamma\,d u_\varepsilon$, which is
$O(\varepsilon)$ since $du$ is bounded.  No derivative of $du$ is
required in this estimate.

Define the stationary smooth spacetime metric from
$u_\varepsilon,X_\varepsilon,g$ as in
Lemma~\ref{lem:variation:stationary-action}.  Its induced second form is
$b_\varepsilon$; let $\mu_\varepsilon,J_\varepsilon$ be its constraint
components.  Since $b_\varepsilon\to K$ in $C^1$,
\[
 \mu_\varepsilon\to\mu,\qquad J_\varepsilon\to J,\qquad
 S_{b_\varepsilon}(u_\varepsilon,X_\varepsilon)\to S
\]
uniformly.  The smooth action identity and the convolved Hessian
equation consequently give
\[
 u_\varepsilon\mathbf G^\varepsilon_{\mathrm{spatial}}
   =P_\varepsilon-J_{(i}X_{j)}+o_{C^0}(1),\qquad
 P_\varepsilon=(\operatorname{tr}_{g(x)}M_\varepsilon)g(x)
                                                     -M_\varepsilon.
\]
Here the regular terms cancel because
$B-(\operatorname{tr}_gB)g+S=0$, and the volume term tends to zero by
$u\mu+J(X)=0$.  The matrix $P_\varepsilon$ is positive semidefinite:
in a $g(x)$-orthonormal eigenbasis for $M_\varepsilon$ its eigenvalues
are the sums of all but one of the nonnegative eigenvalues of
$M_\varepsilon$.  It is essential that this trace reversal be applied
\emph{after} convolution, at $g(x)$.  There is no estimate here
commuting a variable trace-reversal coefficient past an unbounded
Hessian.  For the area measure in the statement, the corresponding
unmollified positive spatial measure is exactly
\[
 \frac c2\int(g-\nu_D^\flat\otimes\nu_D^\flat)\,dA_D\,d\pi.
\]

Let $\mathbf n_\varepsilon$ be the future slice normal.  For the
$g$-unit normal $\nu$ of $\Sigma$, Gauss--Codazzi and nullness yield
\[
 \begin{split}
 u_\varepsilon\mathbf{Ric}^\varepsilon
       (\mathbf n_\varepsilon\pm\nu,
        \mathbf n_\varepsilon\pm\nu)
 &=u_\varepsilon\mathbf G^\varepsilon
       (\mathbf n_\varepsilon\pm\nu,
        \mathbf n_\varepsilon\pm\nu)\\
 &\geq u\mu\pm2uJ(\nu)-J(\nu)g(X,\nu)-o(1)\geq-o(1).
 \end{split}
\]
To check the last inequality pointwise, if $\mu=0$ then $J=0$.
Otherwise complementarity and $u\geq|X|$ imply
$|J|=\mu$, $|X|=u$, and $J=-\mu X^\flat/u$.  The expression is
therefore $u\mu(1\pm J(\nu)/\mu)^2$.  This verifies the null Ricci
lower bound uniformly, without bounding $M_\varepsilon$ or the full
spacetime curvature.

We give the smooth focusing identity used to finish.  Write
$\theta_\varepsilon=H_\Sigma+\operatorname{tr}_\Sigma b_\varepsilon$
and $\ell_\varepsilon=\mathbf n_\varepsilon+\nu$.  Launch orthogonal
null geodesics from the patch, and intersect their local null
hypersurface with the stationary slices.  For each fixed
$\varepsilon$ this is a smooth construction for some positive time;
no uniform time interval is needed.  On each slice normalize the
null generator by $\ell=\mathbf n+\nu$.  If $\chi$ denotes its screen
second form, commuting covariant derivatives along an affinely
parametrized generator in a parallel screen frame gives
\[
 \nabla_\ell\chi_{ab}=-\chi_{ac}\chi_{cb}-\mathcal R_{ab},
 \qquad \sum_a\mathcal R_{aa}=\mathbf{Ric}(\ell,\ell),
\]
where $\mathcal R$ is the screen Jacobi curvature matrix.
Tracing gives $\ell(\theta)=-|\chi|^2-\mathbf{Ric}(\ell,\ell)$ for
that affine normalization.  For the slice normalization,
$\boldsymbol\nabla_\ell\ell=\kappa\ell$, so the trace identity instead
has the additional term $\kappa\theta$.  Differentiating
$\mathbf g(\ell,\mathbf n)=-1$ gives
\[
 \kappa=\mathbf g(\ell,\boldsymbol\nabla_\ell\mathbf n)
          =u^{-1}\nu(u)+b(\nu,\nu),
\]
because the slice-normal acceleration is $\nabla\log u$.

The velocity with slice coordinate $z$ as parameter is $u\ell$.
Subtracting the stationary Killing translation
$\partial_z=u\mathbf n+X$ leaves the spatial velocity $u\nu-X$,
whose normal component is $Q=u-X\cdot\nu$ and tangential component
is $-X^T$.  At the initial slice $\chi=A_\Sigma+b|_{T\Sigma}$.
Thus the two descriptions of the expansion derivative give exactly
\[
 \begin{split}
 L_{+,\varepsilon}Q_{+,\varepsilon}
       -X_\varepsilon^T(\theta_\varepsilon)
  ={}&-u_\varepsilon\bigl(
       |A_\Sigma+(b_\varepsilon)_{\mathrm{tan}}|^2
       +\mathbf{Ric}^\varepsilon(\ell_\varepsilon,
                                  \ell_\varepsilon)\bigr)\\
    &+\bigl(\nu(u_\varepsilon)
              +u_\varepsilon b_\varepsilon(\nu,\nu)\bigr)
                                                    \theta_\varepsilon.
 \end{split}
\]
The notation $L_{+,\varepsilon}$ means the displayed spatial
linearization with $K$ replaced by $b_\varepsilon$, without assuming
that $\theta_\varepsilon$ is zero.

Since $H_\Sigma=\operatorname{tr}_\Sigma K=0$ and
$b_\varepsilon\to K$ in $C^1$, one has
$\theta_\varepsilon\to0$ in $C^1$ on the patch.  Both the tangential
transport term and the final normalization term therefore tend
uniformly to zero; their coefficients involve only bounded first
jets.  The null Ricci lower bound gives
\[
 L_{+,\varepsilon}Q_{+,\varepsilon}
       \leq-u_\varepsilon
                 |A_\Sigma+(b_\varepsilon)_{\mathrm{tan}}|^2+o(1).
\]
The principal part of all these operators is the fixed
$-\Delta_\Sigma$.  Their first-order and zeroth-order coefficients
converge uniformly, while $Q_{+,\varepsilon}\to Q_+$ uniformly with
uniform local Lipschitz bounds.  Pairing with a compactly supported
smooth test function and integrating the Laplacian by parts thus
passes to the asserted distributional inequality.  For the drift
term, the gradients converge weakly against integrable test fields,
and the error in its coefficient tends uniformly to zero.

Finally replace stationary time $z$ by $-z$ and the future normal by
its negative.  This replaces $(K,X,J)$ by $(-K,-X,-J)$ while leaving
$u,g,S,B,\mathcal M$ unchanged.  Applying the plus calculation gives
$Q_-=u+X\cdot\nu$ and the claimed minus inequality.  The argument is
local on every compact subset of $\{u>0\}$, so a partition of unity
establishes the statement for all test functions specified in the
lemma.
\end{proof}

\subsection{Capacity estimates and positive Jacobi fields at singular ends}

The maximum-principle argument will be applied to $Q_\pm/\varphi$.
To ensure that a positive maximum cannot escape into the singular
set, we prove that $1/\varphi$ is bounded and tends to zero there.
Capacity cutoffs allow bounded subsolutions to cross the singular
set; dilation and a cone Liouville theorem then give the required
vanishing along every singular approach.

We first specify the geometric measure theory used in this subsection.
Let $\Gamma$ be the compact support of the perimeter measure of a set
whose boundary is locally perimeter minimizing in an open subset of a
smooth $(k+1)$-dimensional Riemannian manifold.  The open subset contains
$\Gamma$; thus there is no boundary condition or obstacle in this
subsection.  Write $\mathcal S=\operatorname{sing}\Gamma$ and
$d\mu=d\mathcal H^k_g\lfloor\Gamma$.
The interior regularity, monotonicity, and compactness theorems for
minimizing boundaries give the following facts \cite{Simon2018}:
$\Gamma\setminus\mathcal S$ is a smooth embedded minimal hypersurface,
$\mathcal S$ is closed with Hausdorff dimension at most $k-7$, and
$\mu(B_r(x))\leq\Lambda r^k$ for $x\in\Gamma$ and all sufficiently small
$r$, with uniform constants on the compact ambient neighborhood.
The singular set is empty when $k\leq6$.
Dilations about a fixed point have subsequences converging to
multiplicity-one perimeter-minimizing boundary cones, with convergence
of perimeter measures and supports on compact sets, and smooth
multiplicity-one graphical convergence on compact subsets of the
regular limit.  A multiplicity-one plane tangent cone implies
regularity of the original point.  These statements concern boundaries
of sets; no multiplicity assumption is being inferred merely from
stationarity of a varifold.

\begin{lemma}[Removal of the singular set for bounded subsolutions]
\label{lem:variation:capacity-submean}
Suppose $k>2$ and let $\Gamma$ be as above.  A bounded nonnegative
function $w\in W^{1,2}_{\mathrm{loc}}(\Gamma\setminus\mathcal S)$
satisfying $\Delta_\Gamma w\geq-F$ weakly there, where $F\geq0$ is
constant, has locally finite Dirichlet energy through $\mathcal S$.
The inequality extends to compactly supported tests on all of
$\Gamma$.  There are $r_0>0$ and $C_0<\infty$, depending only on the
ambient neighborhood, $k$, and $\Lambda$, such that
\[
 \mathop{\rm ess\,sup}_{\Gamma\cap B_{r/2}(x)} w
 \leq C_0\left[
       \left(r^{-k}\int_{\Gamma\cap B_r(x)}w^2\,d\mu\right)^{1/2}
       +Fr^2\right],\qquad 0<r\leq r_0.
\]
For a continuous function on the regular part the essential supremum
can be replaced by the supremum there.  The constants are uniform
under dilations by scales tending to zero about points in the fixed
compact ambient neighborhood.  One may in particular start with $w$
on one regular component and extend it by zero to the other regular
components.
\end{lemma}

\begin{proof}
Here is the capacity calculation and the energy argument needed before
using a Sobolev inequality.  On a fixed compact region cover its
singular set by finitely many ambient balls $B_{r_i}(p_i)$, with
$r_i<\varepsilon$ and $\sum_i r_i^{k-2}<\varepsilon$.
Such covers exist because
$\dim_{\mathcal H}\mathcal S\leq k-7<k-2$.
Let $\theta_i$ be one on $B_{r_i}(p_i)$, zero outside $B_{2r_i}(p_i)$,
and satisfy $|\nabla\theta_i|\leq2/r_i$.
The cutoff $\chi=1-\max_i\theta_i$ satisfies
\[
 \int_\Gamma|\nabla\chi|^2\,d\mu
 \leq C\Lambda\sum_i r_i^{k-2},\qquad
 \int_\Gamma(1-\chi)^2\,d\mu
 \leq C\Lambda\sum_i r_i^k.
\]
The first estimate follows pointwise almost everywhere by selecting an
index attaining the maximum.  Taking shrinking covers gives
$0\leq\chi_j\leq1$, equal to zero near $\mathcal S$, converging to one
at every regular point, with both displayed integrals tending to zero.
Local versions use a slightly larger compact region than the support
of the test under consideration.

For a fixed smooth ambient cutoff $\eta$, test the subsolution
inequality with $\eta^2\chi_j^2w$.  Sobolev approximation on the smooth
regular part justifies this test.  Expanding and applying Young's
inequality gives
\[
 \int \eta^2\chi_j^2|\nabla w|^2\,d\mu
 \leq 4\int w^2|\nabla(\eta\chi_j)|^2\,d\mu
      +2F\int\eta^2\chi_j^2w\,d\mu.
\]
The right side is uniformly bounded because $w$ is bounded, the area
is locally finite, and the capacity energies tend to zero.  Fatou's
lemma proves local finite energy.  Moreover, $\chi_jw$ converges to
$w$ in local $W^{1,2}$: the extra gradient $w\nabla\chi_j$ tends to
zero in $L^2$, and dominated convergence applies to
$(1-\chi_j)\nabla w$.
For a nonnegative smooth compact test $\psi$, the tests
$\psi\chi_j^2$ now pass to the limit in the weak inequality.  Indeed,
the only extra term is bounded by
\[
 2\|\psi\|_\infty
 \|\nabla w\|_{L^2(\operatorname{supp}\psi)}
 \|\nabla\chi_j\|_{L^2},
\]
and tends to zero.  This proves removability, without assuming energy
regularity in advance.

The required Sobolev input is the Michael--Simon inequality
\cite{MichaelSimon1973}, in the following local $L^2$ form:
\[
 \left(\int|v|^{2k/(k-2)}\,d\mu\right)^{(k-2)/k}
 \leq C\int\bigl(|\nabla v|^2+C_{\mathrm{amb}}v^2\bigr)\,d\mu.
\]
Initially $v$ is compactly supported in the regular part.  To obtain
this form, apply the $L^1$ Michael--Simon inequality to
$|v|^{2(k-1)/(k-2)}$ and use Cauchy--Schwarz.  A smooth isometric
embedding of a slightly larger compact ambient neighborhood gives a
bound for its second fundamental form.  Since $\Gamma$ is minimal
in the ambient metric, its Euclidean mean curvature in that embedding
is bounded by the trace of this ambient second fundamental form.
In particular the bound does not involve $|A_\Gamma|$.
The capacity approximation just proved extends the inequality to the
tests used here.  Dilating this fixed embedding multiplies its
second fundamental form by the dilation scale, which also proves the
claimed uniformity in blow-ups.

For completeness, put $v=w+Fr^2$; if $F=0$, first put $v=w+\delta$
and later let $\delta\downarrow0$.  On $B_r(x)$ this is a positive
subsolution of $\Delta v\geq-r^{-2}v$.
Testing with $\eta^2v^{p-1}$, $p\geq2$, and applying the preceding
Sobolev inequality gives
\[
 \left(\int|\eta v^{p/2}|^{2k/(k-2)}\,d\mu\right)^{(k-2)/k}
 \leq Cp^2\int
       (|\nabla\eta|^2+r^{-2}\eta^2)v^p\,d\mu,
\]
after decreasing $r_0$ so the bounded ambient term is absorbed in
$r^{-2}$.  All power tests follow from boundedness of $w$ and the
energy approximation above.  Iteration with
$p_j=2(k/(k-2))^j$ on radii
$r/2+r/2^{j+1}$ yields the displayed submean estimate; the factors
converge since $\sum_j(1+j)/p_j<\infty$.
The contribution of the added constant is bounded by $CFr^2$ using
$\mu(B_r)\leq\Lambda r^k$.
Finally, every regular point has a connected smooth neighborhood in
one regular component.  The boundary within $\Gamma$ of any regular
component is contained in $\mathcal S$.  Zero extension to other
components therefore adds no regular interface, and the same
capacity argument applies.
\end{proof}

\begin{lemma}[A flat link component determines the entire cone]
\label{lem:variation:flat-link-component}
Let $C\subset\mathbb R^{k+1}$ be the support of a nonzero
multiplicity-one perimeter-minimizing boundary cone.  If a component
of the regular link $L_{\mathrm{reg}}=(C\cap S^k)_{\mathrm{reg}}$
is totally geodesic in $S^k$, then $C$ is a hyperplane.
\end{lemma}

\begin{proof}
Write $d=k-1$ for the link dimension and
$\mathcal S_L=(\operatorname{sing}C)\cap S^k$.
The product description of the cone away from its vertex and the
singular-dimension estimate give
$\dim_{\mathcal H}\mathcal S_L\leq k-8=d-7$ when the link is
singular.  Its area growth is of order $r^d$: on a fixed annulus the
radial product formula and cone area growth imply this bound by
integrating over radial intervals of length comparable to $r$.
Thus $\mathcal S_L$ has zero $2$-capacity in the link by the preceding
cover construction.  A smooth link needs no capacity cutoffs.

Let $U$ be the totally geodesic regular component.  It lies in a
great equator $E$ and is open in $E$.  Every limit point of $U$ that
is regular belongs to $U$: a connected local regular chart through
that point meets $U$, so belongs to the same component.  Consequently
$U$ is relatively closed as well as open in $E\setminus\mathcal S_L$.
This latter set is connected.  One elementary justification is to
use stereographic coordinates on $E$.  Given two points outside
$\mathcal S_L$, almost every intermediate point can be joined to both
by straight segments missing $\mathcal S_L$: the set of bad
intermediate points is contained in the two radial shadows of that
set, whose Hausdorff dimension is at most
$\dim_{\mathcal H}\mathcal S_L+1<d$.  Taking countably many bounded
pieces handles the point omitted by stereographic projection.
Thus $U=E\setminus\mathcal S_L$, and closedness of the support implies
$E\subset C\cap S^k$.

Any other connected regular link component $V$ is disjoint from $E$.
Indeed a smooth regular neighborhood of a point of $E$ already
contains the equatorial sheet; as an embedded hypersurface of the
same dimension it is that sheet locally.  By connectedness $V$ lies
in one strict hemisphere.  Its suitably signed equatorial height
$h$ satisfies
\[
 h>0,\qquad |h|\leq1,\qquad |\nabla_Vh|\leq1,
 \qquad \Delta_Vh=-d h.
\]
This Laplace identity holds for coordinate functions on a minimal
$d$-dimensional submanifold of the unit sphere.
Test it on $V$ with the squared link capacity cutoffs $\chi_j^2$.
These tests have compact support in $V$: any relative boundary of
$V$ in the compact link is singular.  They give
\[
 d\int_Vh\chi_j^2\,d\mathcal H^d
 =2\int_V\chi_j\langle\nabla h,\nabla\chi_j\rangle\,d\mathcal H^d
 \longrightarrow0.
\]
Here Cauchy--Schwarz, the finite link area, and the vanishing
capacity energy justify the limit.  Dominated convergence gives
$\int_Vh=0$, contradicting $h>0$ on a nonempty open component.
Hence no such $V$ exists.  The regular support is contained in $E$;
the singular set has zero link area, and the regular set is dense in
the support.  Thus the support is precisely $E$, and multiplicity
one makes the cone a hyperplane.
\end{proof}

\begin{lemma}[A Liouville theorem for minimizing cones]
\label{lem:variation:cone-liouville}
Let $C$ be a nonflat multiplicity-one perimeter-minimizing boundary
cone of dimension $k$.  If $z$ is a bounded, nonnegative, smooth
function on $\operatorname{reg}C$ satisfying
\[
 \Delta_C z\geq |A_C|^2z,
\]
then $z=0$.  No smoothness or connectedness of the regular link is
assumed.
\end{lemma}

\begin{proof}
A nonflat minimizing boundary cone has $k\geq7$.
In logarithmic radius $s=\log r$ and the regular link $L$, set
$a=k-2>0$, $V=|A_L|^2$, and $Z(s,\omega)=z(e^s\omega)$.
The cone metric and the scaling of the second fundamental form give
\[
 Z_{ss}+aZ_s+\Delta_L Z\geq VZ.
\]
The operator on the left is in divergence form for the measure
$e^{as}\,ds\,d\mathcal H^{k-1}_L$.
On any compact $s$-strip, test after dropping $VZ\geq0$ with
$\eta(s)^2\chi_j(\omega)^2 Z$.
The energy proof in Lemma~\ref{lem:variation:capacity-submean}, with
this smooth positive weight, proves finite energy on smaller strips
and permits removal of the link capacity cutoffs.  In particular,
tests constant in the link variable are legitimate.  It follows that
\[
 q(s)=\int_LZ(s,\omega)\,d\mathcal H^{k-1}_L(\omega)
 \quad\hbox{satisfies}\quad q''+a q'\geq0
\]
on the entire real line in distributions.  The energy bound gives
$q\in H^1_{\mathrm{loc}}(\mathbb R)$, so we use its continuous
representative.  It is bounded and nonnegative.

Such a $q$ is nondecreasing.  To see this, convolve with a smooth
nonnegative kernel.  For the resulting smooth bounded function,
$(e^{as}q')'\geq0$.  If $q'(s_0)<0$, then for every $s<s_0$
\[
 q'(s)\leq e^{a(s_0-s)}q'(s_0).
\]
Integrating backwards forces $q(s)\to+\infty$ as $s\to-\infty$,
contrary to boundedness.  Thus all mollifications are nondecreasing,
and so is $q$.

Suppose $Z$ is not zero.  Positivity on a regular open patch implies
$q(s_0)>0$ for some $s_0$.  Choose
$\rho\in C_c^\infty((1,2))$, $\rho\geq0$, $\int\rho=1$, and set
\[
 \rho_T(s)=T^{-1}\rho(s/T),\qquad
 Z_T(\omega)=\int_{\mathbb R}\rho_T(s)Z(s,\omega)\,ds.
\]
For $T>s_0$ one has $\int_L Z_T\geq q(s_0)>0$, while
$0\leq Z_T\leq\|z\|_\infty$.  On every compact regular link patch,
integration by parts in $s$ gives
\[
 \Delta_LZ_T\geq VZ_T-e_T,
 \qquad
 |e_T|\leq\|z\|_\infty
       (\|\rho_T''\|_{L^1}+a\|\rho_T'\|_{L^1})
       \leq C/T.
\]
Local Caccioppoli estimates give uniform $W^{1,2}$ bounds there.
Weak compactness on a countable exhaustion by regular patches and
weak-star compactness in $L^\infty(L)$ provide a common subsequence with limit
$\overline Z\in W^{1,2}_{\mathrm{loc}}(L)$ satisfying
\[
 0\leq\overline Z\leq\|z\|_\infty,\qquad
 \int_L\overline Z\geq q(s_0),\qquad
 \Delta_L\overline Z\geq V\overline Z.
\]
The integral passes to the weak-star limit since the whole link has
finite area.  There is no assumption here that $V$ is bounded or
integrable at its singular set.

Test the last inequality with $\chi_j^2\overline Z$, supported away
from the singular link.  Expanding and using Young's inequality gives
\[
 \frac12\int_L\chi_j^2|\nabla\overline Z|^2
 +\int_LV\chi_j^2\overline Z^2
 \leq2\|z\|_\infty^2\int_L|\nabla\chi_j|^2
 \longrightarrow0.
\]
Fatou's lemma shows that $\overline Z$ is constant on each regular
link component, and that any component with positive constant has
$A_L=0$.  There are at most countably many regular components, since
the regular link is a second countable manifold.  Its positive
integral therefore supplies at least one such component.
Lemma~\ref{lem:variation:flat-link-component} would make $C$ a plane,
contrary to hypothesis.  Hence $Z$, and therefore $z$, is zero.
\end{proof}

The cone Liouville theorem supplies the obstruction at a singular
limit. We now apply it to the reciprocal Jacobi field, using the
uniform submean estimate to control approaches that do not remain
in regular cone patches.

\begin{lemma}[Reciprocal Jacobi fields under dilation]
\label{lem:variation:reciprocal-compactness}
Let $\Sigma$ be one connected regular component of $\Gamma$, and
suppose $\varphi>0$ is smooth on $\Sigma$ with
\[
 \Delta_\Sigma\varphi+
 (|A_\Sigma|^2+\operatorname{Ric}_g(\nu,\nu))\varphi=0,
 \qquad \varphi\geq m_0>0.
\]
Extend $z=1/\varphi$ by zero on the other regular components.
Fix $p\in\Gamma$ and $r_j\downarrow0$.  Pass to a subsequence on
which dilation about $p$ by $r_j^{-1}$ converges to a tangent cone
$C$.  After a further subsequence, the dilated $z_j$ converge smoothly
on all compact regular cone patches to a bounded nonnegative smooth
$z_0$, with
\[
 \Delta_Cz_0\geq |A_C|^2z_0.
\]
Different regular cone components are allowed to have zero limits.
\end{lemma}

\begin{proof}
In exponential coordinates at the fixed point $p$, let $g_j$ be the
rescaled metrics; these tend smoothly to the Euclidean metric on
compact sets.  On the rescaled copy of $\Sigma$,
\[
 \Delta_{\Gamma_j}z_j
 =(|A_{\Gamma_j}|^2+
       \operatorname{Ric}_{g_j}(\nu_j,\nu_j))z_j
       +2|\nabla z_j|^2/z_j
 \geq(|A_{\Gamma_j}|^2-Cr_j^2)z_j.
\]
The zero extension is bounded by $m_0^{-1}$ and is smooth on every
regular component.  On a connected regular limit chart, smooth
multiplicity-one convergence identifies $\Gamma_j$ with a single
connected graph.  The graph lies either in the rescaled $\Sigma$ or
in another component.  Passing to a subsequence fixes this alternative.
In the latter case $z_j=0$ there.  In the former choose a base point
in the chart and normalize $\varphi_j$ by its value there.
The scalar Harnack inequality and interior Schauder estimates
\cite{GilbargTrudinger2001} apply: the graph metrics are smoothly
uniformly elliptic and the Jacobi potentials are uniformly bounded
with every derivative on a slightly larger regular chart.
They bound the normalized positive functions and their reciprocals
in every interior $C^l$ norm.  Therefore $z_j$ has uniform interior
$C^l$ bounds, even if its value at the base point tends to zero; in
that case the same normalization proves smooth convergence to zero.
If its value has a positive limit, reciprocal convergence proves
the claimed differential inequality.  It is also valid on zero
limit patches.  A countable exhaustion by relatively compact
regular charts and diagonal extraction produce compatible limits
on the entire regular cone.
\end{proof}

\begin{lemma}[Positive lower bound and blow-up at every singular end]
\label{lem:variation:jacobi-singular-ends}
Let $\Sigma$ be a connected regular component of the compact locally
perimeter-minimizing boundary $\Gamma$.  Every smooth positive
solution of
\[
 \Delta_\Sigma\varphi+
 (|A_\Sigma|^2+\operatorname{Ric}_g(\nu,\nu))\varphi=0
\]
satisfies $\inf_\Sigma\varphi>0$.  Moreover,
$\varphi(x_j)\to+\infty$ for every sequence $x_j\in\Sigma$ with
$\operatorname{dist}(x_j,\mathcal S)\to0$.
\end{lemma}

\begin{proof}
If $\Gamma$ is smooth, its regular components are compact and the
first assertion follows from positivity and continuity; the second
is vacuous.  Otherwise $k\geq7$, so all preceding lemmas apply.
Choose $C\geq\sup|\operatorname{Ric}_g|$ on the compact ambient
neighborhood.  Convex truncation of the Jacobi equation shows that
\[
 \begin{split}
 w_\delta&=(1-\varphi/\delta)_+,\qquad 0\leq w_\delta\leq1,\\
 \Delta_\Sigma w_\delta
 &\geq (|A_\Sigma|^2+\operatorname{Ric}_g(\nu,\nu))
       (\varphi/\delta)\,1_{\{\varphi<\delta\}}
 \geq-C.
 \end{split}
\]
The possible measure at the truncation level has the favorable
nonnegative sign.  Extend these functions by zero to other regular
components and use Lemma~\ref{lem:variation:capacity-submean}.
Fix $r>0$ small enough that $C_0Cr^2<1/4$ and cover $\Gamma$ by
finitely many balls of radius $r/2$ with corresponding $r$-balls in
the allowed ambient neighborhood.  Since $w_\delta\to0$ at every
regular point and the whole perimeter is finite, dominated
convergence makes all the finitely many integral terms less than
$1/4$ for small enough $\delta>0$.  Thus $w_\delta\leq1/2$
everywhere on $\Sigma$, proving $\varphi\geq\delta/2$.
The radius is fixed before $\delta$ is sent to zero.

Set $z=1/\varphi$, extended by zero as before, and write
$M=\|z\|_\infty<\infty$.  It satisfies
$\Delta z\geq(|A_\Sigma|^2-C)z$ on $\Sigma$ and hence
$\Delta z\geq-CM$ on all regular components.
Suppose the second assertion fails.  After passing to a subsequence,
compactness gives a \emph{fixed} $p\in\mathcal S$ with
$x_j\to p$ and $z(x_j)\geq\epsilon>0$.
Set $r_j=4\operatorname{dist}_g(x_j,p)$ and dilate about $p$.
The points $x_j$ then lie in the rescaled $B_{1/2}$.
A tangent-cone subsequence converges to a nonflat cone, since a
plane tangent cone would imply regularity at $p$.
Lemmas~\ref{lem:variation:reciprocal-compactness} and
\ref{lem:variation:cone-liouville} give $z_j\to0$ smoothly on every
compact regular patch of that cone.

This convergence and the uniform bound $M$ imply
\[
 \int_{\Gamma_j\cap B_1}z_j^2\,d\mu_j\longrightarrow0.
\]
To verify the assertion at singular points, choose a neighborhood of
the cone singular set in $\overline B_{3/2}$ whose closure has
arbitrarily small cone measure.  Such neighborhoods exist because
the singular set has zero cone measure; they may be chosen with
measure-zero boundary.  Perimeter-measure convergence bounds the
$\mu_j$-measure of these neighborhoods by the same arbitrarily small
quantity, up to an error tending to zero.  On their complement in
$\overline B_1$, support convergence and finitely many regular
graph charts give uniform convergence of $z_j$ to zero.
The integral on the omitted neighborhood is bounded by $M^2$ times
its measure.  Let first $j\to\infty$ and then the omitted measure
tend to zero.  The cone has zero $k$-dimensional measure on every
sphere of positive radius, so the displayed ball causes no boundary
measure issue.

The rescaled inequality is $\Delta z_j\geq-CMr_j^2$.
The uniform scaled form of
Lemma~\ref{lem:variation:capacity-submean}, on a fixed ball, now
implies $\sup_{\Gamma_j\cap B_{1/2}}z_j\to0$, contradicting
$z(x_j)\geq\epsilon$.  This proves the blow-up along every singular
approach.  In particular, there has been no use of a cone limit for
dilations centered at moving singular points, and no uniqueness of
the tangent cone is needed.
\end{proof}

\subsection{Diffuse second derivatives on the zero set}

The remaining obstruction is a nonatomic family of sheets on which
$u$ vanishes identically. On a transverse line such a family gives
a diffuse second-derivative measure, but the diffuse part of the
second derivative of a nonnegative function cannot charge its
minimum level.

\begin{lemma}[The diffuse BV derivative at a minimum level]
\label{lem:variation:diffuse-bv-minima}
Let $f\geq0$ be locally Lipschitz on an interval $I$, and assume its
almost-everywhere derivative $v=f'$ belongs to
$BV_{\mathrm{loc}}(I)$.  If $D^dv$ denotes the nonatomic part of
$Dv$, then
\[
 |D^dv|(\{f=0\})=0.
\]
\end{lemma}

\begin{proof}
The one-dimensional BV representative has finite left and right
limits $v(t-)$ and $v(t+)$ at every interior point; the set
$J_v=\{t:v(t-)\ne v(t+)\}$ is countable.  Since
$f(t+h)-f(t)=\int_t^{t+h}v(s)\,ds$, the one-sided derivatives of
$f$ equal these traces.  At a zero of the nonnegative $f$,
\[
 v(t-)\leq0\leq v(t+).
\]
Consequently both traces are zero on $\{f=0\}\setminus J_v$.

We record a direct consequence of the scalar BV chain rule
\cite{AmbrosioFuscoPallara2000}: for every $a\in\mathbb R$,
\[
 |D^dv|(\{t:v(t-)=v(t+)=a\})=0.
\]
It suffices to prove this on a relatively compact interval, where
$|Dv|$ is finite.  Choose $\eta\in C_c^\infty(\mathbb R)$ with
$0\leq\eta\leq1$ and $\eta(0)=1$, and define
\[
 H_\varepsilon(s)=\int_a^s\eta((r-a)/\varepsilon)\,dr.
\]
Then $\|H_\varepsilon\|_\infty\leq C\varepsilon$.
The $C^1$ BV chain rule writes
\[
 D(H_\varepsilon(v))
 =\eta((\widetilde v-a)/\varepsilon)D^dv
 +\sum_{t\in J_v}
       [H_\varepsilon(v(t+))-H_\varepsilon(v(t-))]\delta_t,
\]
where $\widetilde v$ is the common trace off the jump set.
The first term converges in total variation to
$1_{\{\widetilde v=a\}}D^dv$ by dominated convergence with respect
to $|D^dv|$.  Every summand in the second term tends to zero, and
its absolute value is bounded by $|v(t+)-v(t-)|$; the latter jump
sizes have finite sum.  Thus the second term tends to zero in total
variation.  On the other hand $H_\varepsilon(v)\to0$ uniformly, so
its distributional derivative tends to zero.  The restricted signed
measure $1_{\{\widetilde v=a\}}D^dv$ is therefore zero, and hence
so is its total variation.  This proves the asserted level locality.
Apply it with $a=0$ and use that the nonatomic measure $|D^dv|$
does not charge the countable set $J_v$.
Jump atoms are genuinely excluded from the conclusion: $f(t)=|t|$
has $D^2f=2\delta_0$.
\end{proof}

\begin{lemma}[Nonatomic graph mass cannot be supported where the lapse vanishes]
\label{lem:variation:non-atomic-slicing}
Let $V\times I$ be a coordinate cylinder, where
$V\subset\mathbb R^k$ is open, and let
$\{t=h_a(y):a\in A\}$ be a Borel family of uniformly smooth
graphs defined over $V$ in a larger coordinate range.  Only their
portions with $h_a(y)\in I$ are used, so graphs may cross the top or
bottom of the cylinder.  Let $\sigma$ be a finite Borel measure on the parameter
space $A\subset\mathbb R$.  Assume that, for every $y\in V$,
$a\mapsto h_a(y)$ is injective, and that the weights $w(y,a)$ are
measurable with $0<c_0\leq w\leq C_0$ on smaller cylinders.
Suppose a nonnegative locally Lipschitz function $u$ satisfies
\[
 \partial_t^2u=b(y,t)\,dy\,dt+
 \int_Aw(y,a)\,dy\,\delta_{h_a(y)}(dt)\,d\sigma(a)
\]
as distributions on $V\times I$, with each Dirac measure restricted
to $I$, where $b$ is locally bounded and
$b=0$ for Lebesgue-almost every point of $\{u=0\}$.
Write $\sigma=\sigma_{\mathrm{at}}+\sigma_{\mathrm{na}}$ for its
atomic and nonatomic parts.  Then for
$\sigma_{\mathrm{na}}$-almost every $a$,
\[
 |\{y\in V:h_a(y)\in I,\ u(y,h_a(y))=0\}|=0.
\]
In particular $u$ cannot vanish identically on any nonempty open
portion of such a graph inside the cylinder.
The conclusion holds simultaneously in any specified countable
collection of these cylinders, outside one null set of the original
measure on minimizing boundaries from which their parameter
measures are pushed forward.
\end{lemma}

\begin{proof}
For almost every $y$ the equation slices to
\[
 D^2u_y=b(y,t)\,dt+\lambda_y,
 \qquad
 \lambda_y=\int_Aw(y,a)\bigl(\delta_{h_a(y)}|_I\bigr)\,d\sigma(a),
 \qquad u_y(t)=u(y,t).
\]
Here is a precise way to arrange the slice exceptional sets.
Test the cylinder equation first with products of a smooth base
test and each member of a countable collection of smooth tests dense
in $C^2$ on every compact subinterval of $I$.  Fubini gives the equality outside a
single base null set for all tests in that collection.  The local
Lipschitz bound for $u_y$ and the locally finite measure bounds on
the right extend it to every test by density.  They also show that
$u_y'\in BV_{\mathrm{loc}}(I)$.  A countable exhaustion by smaller
cylinders yields one allowable exceptional set for the entire
cylinder.  The vanishing condition on $b$ likewise holds on
$\{u_y=0\}$ for almost every $y$, by Fubini.

For each fixed $y$, injectivity of the graph parameter gives
\[
 \lambda_y(\{t\})
 =\int_{\{a:h_a(y)=t\}}w(y,a)\,d\sigma(a).
\]
Every fiber has at most one member.  Consequently the contribution
of $\sigma_{\mathrm{na}}$ is precisely the nonatomic part
$\lambda_y^{\mathrm{na}}$, and the contribution of
$\sigma_{\mathrm{at}}$ is atomic.  The strictly positive bounded
weights preserve this distinction.
Lemma~\ref{lem:variation:diffuse-bv-minima} applied to $u_y$ shows
that the diffuse part of $D(u_y')$ does not charge
$Z_y=\{u_y=0\}$.  Taking the diffuse parts in the sliced equation
and using $b(y,\cdot)=0$ almost everywhere on $Z_y$ therefore gives
\[
 \lambda_y^{\mathrm{na}}(Z_y)=0
 \quad\hbox{for almost every }y.
\]
Tonelli's theorem now yields
\[
 0=\int_A\int_V
 w(y,a)\,1_{\{h_a(y)\in I,\ u(y,h_a(y))=0\}}
 \,dy\,d\sigma_{\mathrm{na}}(a).
\]
The lower bound on $w$ proves the assertion about almost every
graph, and no uncountable intersection of slice exceptional sets
has been taken.

To make the final assertion explicit, index the cylinders by
$j\in\mathbb N$.  If $\pi$ is the original measure and
$a_j$ is its graph-height map on the measurable subfamily meeting
cylinder $j$, let $N_j$ be the null set of nonatomic parameters
excluded above.  Since $N_j$ contains no atomic parameters,
$\sigma_j(N_j)=0$, and the pushforward identity gives
$\pi(a_j^{-1}(N_j))=0$.  Discard
$\bigcup_{j=1}^\infty a_j^{-1}(N_j)$ together with the previously
specified null family.  For every remaining minimizing boundary,
every nonempty patch of every regular component is represented in
the countable cover.  Whenever its parameter there is nonatomic,
the asserted zero-set exclusion holds.  In particular a regular
component on which $u$ vanishes identically can only have atomic
parameters in this cover.
\end{proof}

\begin{remark}
For the tensor Hessian equation in the variation argument, the
slicing hypothesis has exactly the form above.  In coordinates
$(y,t)$, put $F_a=t-h_a(y)$.  Then
$\nu_a^\flat=dF_a/|dF_a|_g$, and
$dA_a=\sqrt{\det g}\,|dF_a|_g\,dy$ on the graph.
Contracting the Hessian source twice with $\partial_t$, and dividing
the resulting measure by the ambient density
$\sqrt{\det g}$, gives the weight $|dF_a|_g^{-1}$, times the fixed
positive scalar coefficient in that equation.  Uniform
transversality and the smooth metric give the required upper and
lower bounds.  The connection term in
$\partial_t^2u=\nabla^2u(\partial_t,\partial_t)
+\Gamma^i_{tt}\partial_i u$ belongs to the bounded term.
If that term is homogeneous and linear in $u,X,\nabla X,du$, with
$|X|\leq u$, $u$ locally Lipschitz and $X\in C^{1,1}_{\mathrm{loc}}$,
then it vanishes almost everywhere on $\{u=0\}$: one has $u=X=0$
there, and Sobolev zero-set locality gives $du=\nabla X=0$ almost
everywhere there.  Finally, injectivity of the height parameter is
the no-contact property of distinct local minimal graphs, after
coincident graphs have been identified.
\end{remark}

\subsection{Removal of all interior area support}

We can now finish the localization promised at the beginning of the
section. Atomic sheets are totally geodesic. On a nonatomic sheet the
diffuse BV argument ensures that the lapse is positive somewhere;
the reciprocal Jacobi estimate prevents its comparison maximum from
escaping into a singular point. These two cases give one curvature
bound, which makes smooth outermostness applicable.


\begin{theorem}[The area multiplier is supported on the original boundary]
\label{thm:variation:boundary-support}
Under the rigidity hypotheses and equality, the probability measure
in Proposition~\ref{prop:variation:multiplier} satisfies
\[
 \pi\{C:\Gamma_C=S\}=1.
\]
Consequently, on every compact variation and end prototype,
\begin{equation}\label{eq:variation:boundary-multiplier}
 2k\omega\mathcal E'(h,p)-c a'_S(h)
 =\int_{\mathfrak A}\mathfrak C'(h,p;v)\,d\Lambda
       -\int_S\theta'_+(h,p)\,d\zeta,
 \qquad
 a'_S(h)=\frac12\int_S\tr_Sh\,dA.
\end{equation}
In particular the area source in
Equation~\eqref{eq:variation:hessian} vanishes in the open exterior.  Boundary
mass of the constraint measure $\Lambda$ is not included in this
conclusion.
\end{theorem}

\begin{proof}
For $3\leq n\leq7$, the conclusion was proved above using
Proposition~\ref{prop:variation:smooth-localization}. We therefore assume
$n\geq8$.

\smallskip\noindent\textbf{One full-measure family of graph patches.}

Discard the exceptional family in
Lemma~\ref{lem:variation:contact-dichotomy} and the complement of
$\supp\pi$, which has zero measure because $\mathfrak T$ is a compact
metric space. Use the countable cylinders, height maps $a_j$, and
pushforward measures $\lambda_j$ fixed in
Subsection~\ref{subsec:variation:graph-families}; the measures still
come from the original $\pi$, since the discarded family has zero mass.

The $\partial_t^2u$ component of
Equation~\eqref{eq:variation:hessian}, written in each such cylinder, has exactly
the form in Lemma~\ref{lem:variation:non-atomic-slicing}.
Its positive weight is $c/(2|d(t-h_a(y))|_g)$.
The remaining bounded term is homogeneous and linear in
$u,X,\nabla X,du$.  On $\{u=0\}$, causality implies $X=0$, and
Sobolev locality gives $du=\nabla X=0$ almost everywhere.  The bounded
term thus vanishes there almost everywhere.  Apply that lemma on
each cylinder and remove the countable union of its exceptional
families of minimizers.  Every remaining non-atomic graph patch
has $u>0$ somewhere; indeed its zero set has zero graph area.
This discarding is on the original measure space $\mathfrak T$ by
the pushforward identity, not on an uncountable collection of
separately chosen leaves.

\Needspace{3\baselineskip}
\smallskip\noindent\textbf{A curvature bound on each regular component.}

Fix a remaining wholly interior frontier and a connected component
$\Sigma$ of its regular part.  If a local graph parameter is an atom,
then all the coincident minimizers agree along this entire component.
To see this, the coincidence set is open by the regular-sheet
comparison principle; a limit point in $\Sigma$ lies in the other
closed frontier and the same principle makes coincidence open
there as well.  Connectedness makes the set all of $\Sigma$.
Thus the same positive-mass fiber is present in every smaller
patch of this component, and Lemma~\ref{lem:variation:atomic-flatness}
gives $A_\Sigma=0$ everywhere.

Otherwise take a non-atomic patch of $\Sigma$.
Lemma~\ref{lem:variation:common-jacobi} supplies one positive
$\varphi$ on all of $\Sigma$ solving
Equation~\eqref{eq:variation:common-jacobi}.  It is bounded below and diverges
along every singular approach by
Lemma~\ref{lem:variation:jacobi-singular-ends}.
The functions
\[
 Q_+=u-g(X,\nu),\qquad Q_-=u+g(X,\nu)
\]
are nonnegative and bounded on the compact ambient neighborhood.
By the diffuse-zero-set exclusion, $u>0$ somewhere on the chosen
patch, so at least one $Q_\pm$ is positive somewhere.  For that sign,
$Q_\pm/\varphi$ attains a positive maximum on $\Sigma$: it tends
to zero at every singular end, and on a component without singular
ends the component is compact.  All regular frontier components
have relative boundary only in the singular set, so these are the
only possibilities for escape from compact subsets of $\Sigma$.

At a point where this ratio is positive one has $u>0$.
Lemma~\ref{lem:variation:measure-focusing} therefore gives
\[
 L_\pm Q_\pm\le -u|A_\Sigma\pm K_{\rm tan}|^2,
 \qquad L_\pm\varphi=0.
\]
Dividing by $\varphi$ removes the zeroth-order term: if
$L_\pm=-\Delta+b_\pm\cdot\nabla+V_\pm$, the operator on the
ratio is $-\Delta+(b_\pm-2\nabla\log\varphi)\cdot\nabla$.
Its distributional strong maximum principle applies to the locally
Lipschitz ratio.  For example, apply the weak Harnack inequality to
the nonnegative difference between its maximum and the ratio.
The positive maximum set is open in $\Sigma$ by this principle and
closed by continuity; its positivity ensures that the focusing
inequality holds throughout every neighborhood used in this
argument.  Connectedness makes the ratio constant and positive on
all of $\Sigma$.  The displayed inequality then forces
$A_\Sigma\pm K_{\rm tan}=0$ there.  We conclude in both the atomic
and non-atomic cases that
\begin{equation}\label{eq:variation:curvature-bound}
 |A_\Sigma|\le\sup_{\Gamma_C}|K|.
\end{equation}
The bound is independent of the component.  Each regular frontier
has at most countably many components; the preceding single
full-measure family already makes these conclusions simultaneous.

\smallskip\noindent\textbf{Removal of singular points.}

There can now be no singular point on this frontier.  At such a
point take a fixed-center tangent cone.  Smooth convergence on each
regular cone patch and the rescaled bound
$|A_{\rm scaled}|\le r_j\sup|K|\to0$ make every regular cone
component totally geodesic.  Some regular component exists by the
positive density and singular-dimension bound.  Its regular link
component is totally geodesic, so
Lemma~\ref{lem:variation:flat-link-component} makes the cone a
multiplicity-one plane.  Density regularity then makes the original
point regular, a contradiction.

\smallskip\noindent\textbf{Application of smooth outermostness.}

The entire frontier is therefore a compact smooth embedded two-sided
hypersurface, with finitely many components by compactness and local
regularity.  Its filled minimizer has connected exterior containing
the end, as proved in Lemma~\ref{lem:variation:envelope}.  The closure
of that exterior in $\Omega$ is a smooth closed codimension-zero
submanifold with this entire frontier as its intrinsic boundary.
Since the frontier is wholly interior, it is an enclosing cut of
exactly the kind excluded by the outermostness hypothesis.  Equation~\eqref{eq:variation:three-equations} gives $\theta_+=0$ everywhere
with the normal toward the end, a contradiction.

Thus almost every remaining frontier is $S$.  Integrating its fixed
area derivative yields Equation~\eqref{eq:variation:boundary-multiplier}.
The moment regularity and interior equations were obtained before
this removal, so they remain valid with zero interior area source.
\end{proof}
